\chapter{本体描述语言与工具}
\chapterart{ch04}{语义网技术栈：RDF \(\rightarrow\) RDFS \(\rightarrow\) OWL \(\rightarrow\) SPARQL}

\begin{introbox}
\textbf{【导读】}
前两章回答了「本体由什么构成」与「如何从零构建」，本章解决最后一公里：
把概念与方法落到\textbf{机器能读、推理器能算、查询引擎能答}的标准语言上。
我们自底向上走一遍语义网技术栈：RDF 负责陈述事实，RDFS 给事实加上模式，
OWL~2 提供完整的逻辑表达力，SPARQL 负责把知识问出来；
随后以历史互操作视角盘点高级建模模式与工具生态——Protégé、四个编程库与四个主流推理器。
本章的语法片段可以阅读、编辑和比较；涉及“可复算”的主张则统一从 Semantica 包
装载 RDF Dataset 并执行具名场景。Protégé、Jena 与其他工具用于理解标准生态，
不再构成本书的平行执行入口。
\end{introbox}

\chapterbinding{semantica.chapter\_packages.vol1.ch04}{partial}{blocked}
{RDF/RDFS/OWL 来源资产、CQ1--CQ8、SPARQL 集合及开放/封闭世界场景已入包；
RDF Dataset 与 SPARQL 查询由原生 runtime 执行。}
{Manchester Syntax 只无损保留而不解析；本章没有自有 shapes，场景需调用 ch07；
图同构 round-trip oracle 未声明，外部 \texttt{SERVICE} 在离线环境 fail closed。}

\section{语义网技术栈总览}

学习一门语言之前，先看清它所在的「语言家族」。
本体描述语言不是某个实验室的私有发明，而是一组由万维网联盟
（World Wide Web Consortium, W3C）维护了二十多年的开放标准。
理解这组标准各自的定位与相互关系，是本章一切细节的地基。

\subsection{为什么需要标准语言}

回顾历史能看清动机。二十世纪八九十年代，知识表示领域百花齐放：
KL-ONE 家族的框架语言、KIF 交换格式、Ontolingua 本体库……
每套系统自带语法、语义与工具，知识写进去就被锁在里面，
跨系统交换知识的成本高到几乎等于重写。
那段历史沉淀为一条行业共识：知识表示的价值与「能被多少系统读懂」成正比——
语言不统一，知识就只是另一种形式的孤岛数据。
这与关系数据库标准化之前的处境如出一辙：在 SQL 统一江湖之前，
每家厂商的查询语言互不相认，换数据库约等于换全套应用。

语义网（Semantic Web）运动的核心贡献，是把知识表示推上了
\textbf{Web 规模的标准化轨道}：用全球统一的标识符命名万物，
用极简的图模型陈述事实，再分层叠加模式、逻辑与查询能力。
每一层都是公开标准，任何人都可以实现解析器、存储库与推理器——
这正是标准生态中曾经可以用 Protégé 编辑、用 HermiT 推理、
拿 Fuseki 提供查询服务而三者来自不同团队的原因。对本书而言，
这种互操作性用于读取遗留资产与构造 adapter，不改变 Semantica
作为唯一可执行语义正本的地位。

值得一提的是 W3C「推荐标准」（Recommendation）一词的分量。
一份规范要走完工作草案、候选推荐、建议推荐的全流程，
必须拿出多个相互独立的实现并通过公开测试套件互证，才能盖上推荐标准的章。
这套流程速度不快，换来的是罕见的稳定性：2004 年写下的 RDF 文件，
今天的任何合规解析器仍能一字不差地读出——
对生命周期以十年计的知识工程而言，这种「慢」恰恰是美德。

\subsection{W3C 标准族谱}

本章涉及的标准可以排成一张族谱，每个成员管一摊事：

\begin{itemize}
  \item \textbf{RDF}（Resource Description Framework，资源描述框架）——
        \textbf{数据模型}标准：规定知识的原子是三元组、知识的整体是图。
        1999 年发布首个推荐标准，2004 年形成完整的 RDF 1.0 标准族，
        2014 年升级为 RDF 1.1（全面转向 IRI、Turtle 转正为标准语法）；
        引入「可引用三元组」的 RDF 1.2（即社区熟知的 RDF-star 方向）
        在本书写作时仍处于标准化进程中。
  \item \textbf{RDFS}（RDF Schema）——\textbf{轻量模式}标准：
        在 RDF 之上提供类、类层次、属性层次与定义域/值域词汇，
        2004 年随 RDF 1.0 发布，2014 年同步更新为 RDF Schema 1.1。
  \item \textbf{OWL}（Web Ontology Language，Web 本体语言）——
        \textbf{本体逻辑}标准：提供等价、不相交、枚举、基数约束、
        属性特性等完整描述逻辑表达力。OWL 1 于 2004 年发布，
        OWL 2 于 2009 年发布（2012 年出第二版勘误），
        并带来三个面向不同计算需求的子语言（Profile），见 4.5 节。
  \item \textbf{SPARQL}（SPARQL Protocol and RDF Query Language）——
        \textbf{查询与更新}标准：2008 年的 SPARQL 1.0 奠定图模式匹配查询，
        2013 年的 SPARQL 1.1 补齐聚合、子查询、属性路径、联邦查询与更新语言，
        同时规定了 HTTP 协议层——「SPARQL 端点」因此成为跨厂商的服务接口。
  \item \textbf{SHACL}（Shapes Constraint Language，形状约束语言）——
        \textbf{数据校验}标准：2017 年发布，用「形状」描述数据应满足的结构约束，
        以封闭世界的方式产出违规报告。它与 OWL 的推断定位互补，
        本章 4.4 节讨论二者分工，完整用法见第 7 章。
\end{itemize}

这张族谱有清晰的分工逻辑：RDF 管「怎么说事实」，RDFS/OWL 管「怎么定义词汇与概念」，
SPARQL 管「怎么问」，SHACL 管「怎么验」。
五者共享同一个数据模型，因此其标准语义可以在同一个 Semantica package 中组合——
这是技术栈而非单一语言的含义；多个外部工具各自运行则不自动组成可验收执行链。
查任何细节请认准一手来源：每份标准在 W3C 技术报告页都有带版本日期的正式文本，
搜索引擎给出的二手教程可能停留在旧版本（例如还在教 2004 年的 RDF 1.0 写法），
拿不准时回标准原文核对，是这一行的基本功。

\subsection{「分层蛋糕」：技术栈的结构比喻}

语义网社区习惯用一张「分层蛋糕」（layer cake）图描绘这套体系，
自底向上大致是：最底层是 \textbf{IRI 与 Unicode}——保证万物可命名、全球无歧义；
其上是\textbf{语法层}——Turtle、RDF/XML、JSON-LD 等序列化格式；
再上是 \textbf{RDF 数据层}——三元组与图模型；
然后是\textbf{模式与本体层}——RDFS 提供基础词汇、OWL 提供完整逻辑；
\textbf{规则层}（SWRL 等，见第 5 章）与本体层并肩；
\textbf{SPARQL} 作为查询服务横切各层；
顶部的逻辑、证明与信任层至今仍更多是研究愿景。

这个比喻的工程价值在于两点。
第一，\textbf{每一层只依赖下层的语义}：换一种序列化格式不影响图的内容；
受控的存储投影也可以更换 adapter，而不改变 Semantica package 中内容绑定的本体公理。
“可替换”只适用于投影实现，不意味着各层可以各自宣布语义结论。
第二，\textbf{项目可以「吃到第几层就停」}：
不少知识图谱项目只用到 RDF + RDFS + SPARQL 就已产生业务价值，
OWL 推理按需引入——分层架构允许渐进采纳，而不是全有或全无。

一条常见的渐进路线是三步走。第一步只做「实体目录」：
把设备、物料、工单的台账数据转成内容绑定的 RDF Dataset，纳入 Semantica package，
先解决「数出多门」的标识统一问题；若确有远程查询需要，再从已验收 Dataset
生成只读 SPARQL 投影端点；
第二步补上 RDFS 类层次与 domain/range，让查询能沿层次泛化；
第三步在痛点明确的子领域引入 OWL 公理与推理器——
例如设备能力匹配、工艺约束检查。每一步都独立交付价值，也都为下一步铺路。
比起「先建大而全的本体、再谈应用」的瀑布式路线，
这种爬楼梯的走法在真实组织里存活率高得多。

\subsection{标准化对企业意味着什么}

在企业语境下，「用 W3C 标准」不是情怀，而是实打实的风险控制：

\begin{itemize}
  \item \textbf{投影厂商可替换}。所有合规的三元组库——开源的 Jena TDB2、RDF4J，
        商用的 GraphDB、Stardog、Amazon Neptune 等——讲同一种数据语言与查询语言。
        作为外部投影时，迁移可从 Semantica 已绑定 Dataset 重新生成，或由旧库导出
        N-Triples 后作为候选输入重新绑定、校验；SPARQL 的标准化降低 adapter 迁移成本，
        但迁移完成仍须重跑 package 场景并生成 receipt。
        对比之下，属性图阵营长期依赖各家私有查询语言，
        直到 2024 年才迎来 ISO 的 GQL 标准；
        而私有知识库产品一旦停服，知识资产的迁移成本可能是毁灭性的。
  \item \textbf{数据寿命长于软件寿命}。企业知识库的生命周期以十年计，
        远超任何一代存储产品。把知识沉淀在开放标准格式中，
        相当于给资产上了「格式保险」。
  \item \textbf{生态与人才}。标准化带来可积累的技能与大量遗留实现：
        Protégé、HermiT、Fuseki、rdflib 的产物可以经受控 adapter 交换，
        教材与社区答案也都围绕同一套语言展开；进入本书边界后，交换产物必须回到
        Semantica 做内容绑定、重验证与 receipt，不能自由拼成另一条生产执行链。
\end{itemize}

顺带回答一个高频疑问：「知识图谱用属性图数据库（如 Neo4j）不行吗？」
行，而且各有所长——属性图的边可以直接携带属性，图算法生态成熟；
但其数据模型与查询语言长期缺乏跨厂商标准，
模式层也没有 RDFS/OWL 这样自带推理语义的标准词汇。
第 7 章会给出两类\textbf{投影 adapter} 的详细对照与选型判据；本章的立场很简单：
W3C 技术栈是 Semantica package 的交换语言与内容模型，属性图可以承担分析投影；
无论选择哪种外部存储，都不获得独立于 package/runtime/receipt 的语义权威。

\begin{noteblock}
\textbf{与第 2 章的衔接}：第 2 章介绍的描述逻辑（DL）是这套标准的理论内核——
OWL 的每个构造器都对应一个 DL 构造器，OWL 2 DL 整体对应
\(\mathcal{SROIQ}(D)\)。本章读到 4.5 节时，
你会看到那些抽象符号如何一一变成可敲进编辑器的关键字。
\end{noteblock}

本章行文按技术栈自底向上推进：4.2 与 4.3 节讲数据层（模型与序列化），
4.4 与 4.5 节讲模式与逻辑层（RDFS 与 OWL），4.6 节讲建模惯用法，
4.7 节讲查询层（SPARQL），4.8 节讲工具链，4.9 节小结并给出练习。

\section{RDF 数据模型}

\subsection{三元组与图}

RDF 的设计思想极简：\textbf{一切知识都拆成「主语--谓语--宾语」三元组}
（subject--predicate--object triple）。
「3 号车床位于 A 工作单元」写成
\texttt{(mfg:Lathe\_003, mfg:locatedIn, mfg:WorkCell\_A)}；
「1 号车床功率 22 千瓦」写成
\texttt{(mfg:Lathe\_001, mfg:power, "22.0"\^{}\^{}xsd:decimal)}。
主语和宾语是节点，谓语是连接两个节点的\textbf{带标签有向边}——
一组三元组叠在一起，自然构成一张有向标记图，这就是 \textbf{RDF 图}（RDF graph）。
形式上，RDF 图就是三元组的\textbf{集合}：无序、无重复，
同一主语谓语天然允许多个宾语（多值不需要任何特殊设计）。

三元组还有第二种读法：把谓语看作\textbf{二元谓词}，
\texttt{(Lathe\_003, locatedIn, WorkCell\_A)} 即逻辑断言
「locatedIn(Lathe\_003, WorkCell\_A)」。
这正是第 2 章 ABox 断言的标准形态——RDF 图就是一组用 Web
标识符写成的逻辑事实。「图」与「逻辑」两副眼镜随时切换：
谈存储与查询时用图的眼镜，谈推理与语义时用逻辑的眼镜，
4.4 节讨论 RDFS 衍推时你会体会到双视角的便利。

与关系模型对比能看出图模型的脾气。关系表先定列再填行，
模式变更（加列）是伤筋动骨的操作；RDF 图没有「表结构」，
任何人随时可以对任何资源追加新谓语的三元组，模式演化近乎零成本。
更重要的是\textbf{合并自由}：两张 RDF 图的并集仍是合法的 RDF 图，
来自不同系统、不同厂商的数据只要谈论同一批标识符，倒在一起就完成了初步集成——
这是用图做数据集成的根本优势，第 7 章构建知识图谱时将充分兑现。
一个具体画面帮助记忆：制造执行系统导出设备台账的三元组，
质检系统导出检验记录的三元组，
两边都用 \texttt{mfg:Lathe\_001} 指称同一台车床——
两个文件先后装进同一个库，这台车床的全部知识就自动「长」在同一个节点上，
一行对账代码都不用写。标识符的全局性，
把集成成本从「N 平方个点对点接口」降到「N 份数据各自对齐标准」。

\subsection{节点的三种身份：IRI、字面量、空节点}

RDF 图中的节点只有三种身份，规则简单却必须烂熟：

\begin{itemize}
  \item \textbf{IRI}（Internationalized Resource Identifier，国际化资源标识符）——
        资源的全球唯一名字，如 \nolinkurl{http://example.org/manufacturing#Lathe_001}。
        IRI 是 URI 的国际化版本（允许直接使用 Unicode 字符）；
        而日常说的 URL 是「可解引用」的特例。
        RDF 不要求 IRI 必须能在浏览器里打开，但好的实践（Linked Data 惯例）
        鼓励解引用后返回机器可读的描述。
  \item \textbf{字面量}（Literal）——具体数据值：字符串、数字、日期等，
        可带语言标签或数据类型（下一小节展开）。
  \item \textbf{空节点}（Blank Node）——「存在但未命名」的节点，
        相当于自然语言里的「某个东西」（4.2.4 小节专门讨论它的用途与坑）。
\end{itemize}

三种身份在三元组中的合法位置各不相同：
\textbf{主语}可以是 IRI 或空节点；\textbf{谓语}必须是 IRI；
\textbf{宾语}三者皆可。字面量永远只能当宾语——
「22.0 位于 A 单元」这种话在 RDF 里根本写不出来，语法层面就挡住了一类错误。

IRI 的设计值得在项目第一天就定下规矩，因为它一旦发布就近乎不可更改。
三条久经考验的惯例：第一，\textbf{命名空间用稳定域名加版本无关路径}，
版本号放进本体的注解属性而不是每个词项的名字里；
第二，\textbf{井号与斜杠二选一}——小词表用井号命名空间
（整个文件一次取回），大数据集用斜杠命名空间（每个资源可单独解引用）；
第三，\textbf{标识与含义分离}——IRI 是名字不是说明书，
把部门、产线、型号之类的业务含义编码进 IRI，
将来组织调整时会追悔莫及；人读的信息交给 \texttt{rdfs:label} 承载。

完整 IRI 冗长难读，实践中用\textbf{前缀声明}（prefix declaration）缩写：
声明 \texttt{mfg:} 代表 \nolinkurl{http://example.org/manufacturing#} 之后，
\texttt{mfg:Lathe\_001} 就是完整 IRI 的缩写（称为 CURIE 风格写法）。
下面这段开篇代码展示了 ch04 包资产的前缀区与本体声明：

\begin{codebox}
\codeline{@prefix\ rdf:\ \ <http://www.w3.org/1999/02/22-rdf-syntax-ns\#>\ .}
\codeline{@prefix\ rdfs:\ <http://www.w3.org/2000/01/rdf-schema\#>\ .}
\codeline{@prefix\ owl:\ \ <http://www.w3.org/2002/07/owl\#>\ .}
\codeline{@prefix\ xsd:\ \ <http://www.w3.org/2001/XMLSchema\#>\ .}
\codeline{@prefix\ :\ \ \ \ \ <http://example.org/manufacturing\#>\ .}
\codeline{@prefix\ mfg:\ \ <http://example.org/manufacturing\#>\ .}
\codeblank{}
\codeline{\textcolor{cmtgray}{\#\ ==============================}}
\codeline{\textcolor{cmtgray}{\#\ 本体声明\ (Ontology\ Declaration)}}
\codeline{\textcolor{cmtgray}{\#\ ==============================}}
\codeline{<http://example.org/manufacturing>}
\codeline{\hspace*{4\ttcw}a\ owl:Ontology\ ;}
\codeline{\hspace*{4\ttcw}rdfs:label\ "制造领域本体"@zh\ ;}
\codeline{\hspace*{4\ttcw}rdfs:label\ "Manufacturing\ Domain\ Ontology"@en\ ;}
\codeline{\hspace*{4\ttcw}rdfs:comment\ "用于描述制造领域的设备、工艺、产品和工单知识"@zh\ .}
\end{codebox}
\color{black}

\assetsource{ch04}{rdf-turtle-examples}

注意两个细节。其一，\texttt{rdf:}、\texttt{rdfs:}、\texttt{owl:}、\texttt{xsd:}
四个前缀几乎出现在每份 RDF 文件头部，分别指向四个 W3C 标准命名空间——
它们的展开形式值得记住，读旧文件、调试解析错误时都用得上。
其二，\texttt{:} 与 \texttt{mfg:} 被声明为同一个命名空间，
空前缀让本地名字更紧凑，具名前缀让跨文件引用更清晰，教学资产两者并存以便对照。
最后提醒一句：前缀只是\textbf{文档内}的缩写约定，不进图本身——
同一文件换一套前缀名重写，得到的图分毫不差；
跨团队统一前缀用法（比如固定 \texttt{mfg:} 指向企业命名空间）纯属人类间的便利，
公共词表的惯用前缀则早已约定俗成，社区站点 prefix.cc 收录了这些映射。

\subsection{字面量：语言标签与数据类型}

字面量看似平凡，细节却直接影响查询与集成质量。RDF 1.1 把字面量分成三类：

\textbf{语言标签字面量}：\texttt{"制造领域本体"@zh} 与
\texttt{"Manufacturing Domain Ontology"@en}。
标签采用 BCP~47 规范（\texttt{zh}、\texttt{en}、也可细化为 \texttt{zh-Hans} 等），
类型固定为 \texttt{rdf:langString}。
上面代码里同一个本体挂了中英两条 \texttt{rdfs:label}——
这是\textbf{两条不同的三元组}，互不覆盖。
多语言标注从第一天就做，国际化项目会感谢这个习惯；
查询时用 \texttt{FILTER(lang(?name) = "zh")} 取所需语言（4.7 节示范），
若数据里混有 \texttt{zh-Hans} 这类带子标签的写法，
改用 \texttt{langMatches()} 函数做前缀式匹配更稳妥。

\textbf{类型字面量}：\texttt{"22.0"\^{}\^{}xsd:decimal}、
\texttt{"2024-03-01T10:00:00"\^{}\^{}xsd:dateTime}。
双插符后面的 IRI 指明数据类型，常用的有
\texttt{xsd:string}、\texttt{xsd:integer}、\texttt{xsd:decimal}、
\texttt{xsd:double}、\texttt{xsd:boolean}、\texttt{xsd:date}、
\texttt{xsd:dateTime}、\texttt{xsd:anyURI}。
没写类型的裸字符串，在 RDF 1.1 中默认就是 \texttt{xsd:string}。

两个工程提醒。第一，\textbf{词项相等不同于值相等}：
\texttt{"22"\^{}\^{}xsd:integer} 与 \texttt{"22.0"\^{}\^{}xsd:decimal}
是不同词项但数值相等，SPARQL 的 \texttt{=} 按值比较、
\texttt{sameTerm()} 按词项比较，混淆两者是查询结果「莫名缺行」的常见原因。
第二，\textbf{单位不要塞进数值}：\texttt{"22kW"} 这种字符串既不能比较也不能聚合；
包内教学资产的做法是把数值写入内容绑定 Dataset、
单位写进属性的 \texttt{rdfs:comment}（如「单位：kW」），
更严格的方案是引入 QUDT 之类的单位本体显式建模量纲。

数值类型的选择也有讲究：工程量优先用 \texttt{xsd:decimal}
（十进制精确表示，适合尺寸、功率、金额），
\texttt{xsd:double} 留给确属浮点语义的科学计算值——
二进制浮点的舍入误差若出现在「最大加工直径」这类字段上，排查起来格外伤神。
日期一律走 \texttt{xsd:date} 或 \texttt{xsd:dateTime} 而不是自由字符串，
否则区间过滤与排序都要仰仗不可靠的字典序。

\subsection{空节点：用途与坑}

空节点表达「存在某个东西，它有如下性质，但我不给它全局名字」。
合法用途有三类：描述不值得命名的辅助结构
（例如一个地址、一段配置）；承载 RDF 列表的中间节点——
Turtle 里一对圆括号背后，其实是一串以 \texttt{rdf:first}、\texttt{rdf:rest}
串起、以 \texttt{rdf:nil} 收尾的空节点链；
以及 OWL 公理的 Turtle 形态——4.5 节会看到，
\texttt{owl:Restriction} 匿名类在 Turtle 里就是用空节点 \texttt{[ ... ]} 写的。

但空节点的代价常被低估，三处坑必须心中有数：

\begin{itemize}
  \item \textbf{不可全局引用}。空节点标签（如 \texttt{\_:b1}）只在\textbf{本文档内}有效，
        换个文件、换次解析就失效。两份都含空节点的数据合并后，
        无法判断两个空节点是否「同一个东西」，
        增量更新与数据对账因此变得棘手。
  \item \textbf{SPARQL 行为反直觉}。同一查询内，图模式里的空节点相当于
        不进入结果的临时变量；而\textbf{查询结果中出现的空节点标签只是当次结果集的代号}，
        拿它拼第二条查询去「定位同一个节点」是徒劳的——这是新手高频踩坑点。
  \item \textbf{运维成本}。含空节点的图做差异比较（diff）需要图同构判定，
        版本管理、变更审计都因此复杂化。
\end{itemize}

\begin{noteblock}
\textbf{工程建议}：能起名就起名。确实需要消除空节点时，
可用\textbf{斯科伦化}（Skolemization）——把空节点系统性替换为
\texttt{.well-known/genid/} 路径下的内部 IRI，
RDF 1.1 标准为此专门预留了机制。本书第 9 章的实战数据全部使用具名 IRI，
空节点只保留在 OWL 公理的语法内部。
\end{noteblock}

\section{序列化格式巡礼}

RDF 图是抽象数学对象，落到文件里需要\textbf{序列化格式}（serialization format）。
同一张图可以写成多种文本，格式之间无损互转——
本节边讲语法边强调这个「同图异文」的事实。

\subsection{Turtle 语法精讲}

Turtle（Terse RDF Triple Language）是手写与阅读 RDF 的事实标准，
2014 年随 RDF 1.1 转正为 W3C 推荐标准。它的全部语法要点可以一页讲完：

\begin{itemize}
  \item \textbf{前缀与基地址}：\texttt{@prefix mfg: <...> .} 声明缩写；
        与 SPARQL 同源的大写写法 \texttt{PREFIX}（不带句点）同样合法。
  \item \textbf{句点、分号、逗号}：句点结束一条陈述；
        \textbf{分号}复用主语——「同一个东西的下一条属性」；
        \textbf{逗号}复用主语加谓语——「同一属性的下一个值」。
  \item \textbf{关键字 a}：\texttt{rdf:type} 的语法糖，
        \texttt{mfg:Lathe\_001 a mfg:Lathe} 读作「是一台」。
  \item \textbf{方括号}：\texttt{[ ... ]} 内联一个空节点及其属性；
        \textbf{圆括号}：\texttt{( e1 e2 e3 )} 构造 RDF 列表；
        三引号 \texttt{"""..."""} 写多行字符串。
\end{itemize}

两条锦上添花的简写：数字与布尔字面量可以裸写——
\texttt{mfg:power 22.0} 等价于带 \texttt{xsd:decimal} 类型的完整写法，
\texttt{true} 与 \texttt{false} 对应 \texttt{xsd:boolean}；
\texttt{@base} 声明基地址后，相对 IRI（如 \texttt{<\#Lathe\_001>}）按其展开。
注释从井号到行尾——包内教学资产的分节横线全是纯注释，
解析后不留痕迹。

语法虽小，组织风格决定可读性。看包内教学资产的类定义段——
这是一段值得逐行品味的「标准笔迹」：

\begin{codebox}
\codeline{\textcolor{cmtgray}{\#\ ==============================}}
\codeline{\textcolor{cmtgray}{\#\ 类定义\ (Class\ Definitions)\ -\ TBox}}
\codeline{\textcolor{cmtgray}{\#\ ==============================}}
\codeblank{}
\codeline{\textcolor{cmtgray}{\#\ 设备类层次}}
\codeline{mfg:Equipment\ a\ owl:Class\ ;}
\codeline{\hspace*{4\ttcw}rdfs:label\ "设备"@zh\ ;}
\codeline{\hspace*{4\ttcw}rdfs:label\ "Equipment"@en\ ;}
\codeline{\hspace*{4\ttcw}rdfs:comment\ "制造系统中所有物理设备的基类"@zh\ .}
\codeblank{}
\codeline{mfg:ProcessingEquipment\ a\ owl:Class\ ;}
\codeline{\hspace*{4\ttcw}rdfs:subClassOf\ mfg:Equipment\ ;}
\codeline{\hspace*{4\ttcw}rdfs:label\ "加工设备"@zh\ .}
\codeblank{}
\codeline{mfg:CNCMachine\ a\ owl:Class\ ;}
\codeline{\hspace*{4\ttcw}rdfs:subClassOf\ mfg:ProcessingEquipment\ ;}
\codeline{\hspace*{4\ttcw}rdfs:label\ "数控机床"@zh\ .}
\codeblank{}
\codeline{mfg:Lathe\ a\ owl:Class\ ;}
\codeline{\hspace*{4\ttcw}rdfs:subClassOf\ mfg:CNCMachine\ ;}
\codeline{\hspace*{4\ttcw}rdfs:label\ "车床"@zh\ .}
\codeblank{}
\codeline{mfg:MeasurementEquipment\ a\ owl:Class\ ;}
\codeline{\hspace*{4\ttcw}rdfs:subClassOf\ mfg:Equipment\ ;}
\codeline{\hspace*{4\ttcw}rdfs:label\ "测量设备"@zh\ .}
\codeblank{}
\codeline{\textcolor{cmtgray}{\#\ 产品与物料}}
\codeline{mfg:Product\ a\ owl:Class\ ;}
\codeline{\hspace*{4\ttcw}rdfs:label\ "产品"@zh\ .}
\codeblank{}
\codeline{mfg:Material\ a\ owl:Class\ ;}
\codeline{\hspace*{4\ttcw}rdfs:label\ "材料"@zh\ .}
\codeblank{}
\codeline{mfg:Aluminum\ a\ owl:Class\ ;}
\codeline{\hspace*{4\ttcw}rdfs:subClassOf\ mfg:Material\ ;}
\codeline{\hspace*{4\ttcw}rdfs:label\ "铝合金"@zh\ .}
\codeblank{}
\codeline{\textcolor{cmtgray}{\#\ 工艺与工序}}
\codeline{mfg:Process\ a\ owl:Class\ ;}
\codeline{\hspace*{4\ttcw}rdfs:label\ "加工工艺"@zh\ .}
\codeblank{}
\codeline{mfg:Operation\ a\ owl:Class\ ;}
\codeline{\hspace*{4\ttcw}rdfs:label\ "工序"@zh\ .}
\codeblank{}
\codeline{mfg:Turning\ a\ owl:Class\ ;}
\codeline{\hspace*{4\ttcw}rdfs:subClassOf\ mfg:Operation\ ;}
\codeline{\hspace*{4\ttcw}rdfs:label\ "车削工序"@zh\ .}
\codeblank{}
\codeline{\textcolor{cmtgray}{\#\ 工作单元}}
\codeline{mfg:WorkCell\ a\ owl:Class\ ;}
\codeline{\hspace*{4\ttcw}rdfs:label\ "工作单元"@zh\ .}
\codeblank{}
\end{codebox}

\assetsource{ch04}{rdf-turtle-examples}

读这段代码的正确姿势是把分号读成「并且它还……」：
\texttt{mfg:Lathe} 是一个类，\textbf{并且它还}是 \texttt{mfg:CNCMachine} 的子类，
\textbf{并且它还}有中文标签「车床」。五个设备类通过 \texttt{rdfs:subClassOf}
串成四层层次（设备 \(\rightarrow\) 加工设备 \(\rightarrow\) 数控机床 \(\rightarrow\) 车床），
每个类一段、段间空行、注释分节——这套排版习惯让 Turtle 文件像文章一样可读。
团队协作时值得把它写进代码规范：固定前缀顺序、固定缩进宽度、
类按层次分组、属性按对象/数据分区——
风格统一之后，评审者的注意力才能留给真正要紧的语义变化。
另一个值得注意的工程细节：\textbf{Turtle 对行敏感的排版方式对版本管理极友好}，
每条公理的增删在 \texttt{git diff} 里清清楚楚，
这是团队协作维护本体时选择 Turtle 的重要理由。

\subsection{RDF/XML：同图异文的历史样本}

RDF/XML 是 1999 年随 RDF 诞生的\textbf{第一个}标准序列化格式，
在 Turtle 转正之前的十几年里是唯一选择，因此存量本体大量使用它——
很多扩展名为 \texttt{.owl} 的文件，揭开盖子其实是 RDF/XML。
仓库提供了与 Turtle 文件内容对应的 XML 版本，节选对照如下：

\begin{codebox}
\codeline{<?xml\ version="1.0"\ encoding="UTF-8"?>}
\codeline{\textcolor{cmtgray}{<!--}}
\codeline{\hspace*{2\ttcw}\textcolor{cmtgray}{RDF/XML\ 格式示例\ -\ 制造领域本体}}
\codeline{\hspace*{2\ttcw}\textcolor{cmtgray}{RDF/XML\ Format\ Example\ -\ Manufacturing\ Domain\ Ontology}}
\codeblank{}
\codeline{\hspace*{2\ttcw}\textcolor{cmtgray}{RDF/XML\ 是\ RDF\ 的\ XML\ 序列化格式，是最早的标准格式}}
\codeline{\hspace*{2\ttcw}\textcolor{cmtgray}{虽然较为冗长，但在与XML工具集成时有优势}}
\codeline{\textcolor{cmtgray}{-->}}
\codeline{<rdf:RDF}
\codeline{\hspace*{2\ttcw}xmlns:rdf="http://www.w3.org/1999/02/22-rdf-syntax-ns\#"}
\codeline{\hspace*{2\ttcw}xmlns:rdfs="http://www.w3.org/2000/01/rdf-schema\#"}
\codeline{\hspace*{2\ttcw}xmlns:owl="http://www.w3.org/2002/07/owl\#"}
\codeline{\hspace*{2\ttcw}xmlns:xsd="http://www.w3.org/2001/XMLSchema\#"}
\codeline{\hspace*{2\ttcw}xmlns:mfg="http://example.org/manufacturing\#">}
\codeblank{}
\codeline{\hspace*{2\ttcw}\textcolor{cmtgray}{<!--\ 本体声明\ -->}}
\codeline{\hspace*{2\ttcw}<owl:Ontology\ rdf:about="http://example.org/manufacturing">}
\codeline{\hspace*{4\ttcw}<rdfs:label\ xml:lang="zh">制造领域本体</rdfs:label>}
\codeline{\hspace*{4\ttcw}<rdfs:label\ xml:lang="en">Manufacturing\ Domain\ Ontology</rdfs:label>}
\codeline{\hspace*{2\ttcw}</owl:Ontology>}
\codeblank{}
\codeline{\hspace*{2\ttcw}\textcolor{cmtgray}{<!--\ 类定义：设备\ -->}}
\codeline{\hspace*{2\ttcw}<owl:Class\ rdf:about="http://example.org/manufacturing\#Equipment">}
\codeline{\hspace*{4\ttcw}<rdfs:label\ xml:lang="zh">设备</rdfs:label>}
\codeline{\hspace*{4\ttcw}<rdfs:comment\ xml:lang="zh">制造系统中所有物理设备的基类</rdfs:comment>}
\codeline{\hspace*{2\ttcw}</owl:Class>}
\codeblank{}
\codeline{\hspace*{2\ttcw}\textcolor{cmtgray}{<!--\ 类定义：加工设备（子类）\ -->}}
\codeline{\hspace*{2\ttcw}<owl:Class\ rdf:about="http://example.org/manufacturing\#ProcessingEquipment">}
\codeline{\hspace*{4\ttcw}<rdfs:label\ xml:lang="zh">加工设备</rdfs:label>}
\codeline{\hspace*{4\ttcw}<rdfs:subClassOf\ rdf:resource="http://example.org/manufacturing\#Equipment"/>}
\codeline{\hspace*{2\ttcw}</owl:Class>}
\codeblank{}
\codeline{\hspace*{2\ttcw}\textcolor{cmtgray}{<!--\ 类定义：车床\ -->}}
\codeline{\hspace*{2\ttcw}<owl:Class\ rdf:about="http://example.org/manufacturing\#Lathe">}
\codeline{\hspace*{4\ttcw}<rdfs:label\ xml:lang="zh">车床</rdfs:label>}
\codeline{\hspace*{4\ttcw}<rdfs:subClassOf}
\codeline{\hspace*{8\ttcw}rdf:resource="http://example.org/manufacturing\#ProcessingEquipment"/>}
\codeline{\hspace*{2\ttcw}</owl:Class>}
\codeblank{}
\codeline{\hspace*{2\ttcw}\textcolor{cmtgray}{<!--\ 对象属性：canProcess\ -->}}
\codeline{\hspace*{2\ttcw}<owl:ObjectProperty\ rdf:about="http://example.org/manufacturing\#canProcess">}
\codeline{\hspace*{4\ttcw}<rdfs:label\ xml:lang="zh">能加工</rdfs:label>}
\codeline{\hspace*{4\ttcw}<rdfs:domain\ rdf:resource="http://example.org/manufacturing\#Equipment"/>}
\codeline{\hspace*{4\ttcw}<rdfs:range\ rdf:resource="http://example.org/manufacturing\#Material"/>}
\codeline{\hspace*{2\ttcw}</owl:ObjectProperty>}
\codeblank{}
\codeline{\hspace*{2\ttcw}\textcolor{cmtgray}{<!--\ 数据属性：serialNumber\ -->}}
\codeline{\hspace*{2\ttcw}<owl:DatatypeProperty\ rdf:about="http://example.org/manufacturing\#serialNumber">}
\codeline{\hspace*{4\ttcw}<rdfs:label\ xml:lang="zh">序列号</rdfs:label>}
\codeline{\hspace*{4\ttcw}<rdfs:domain\ rdf:resource="http://example.org/manufacturing\#Equipment"/>}
\codeline{\hspace*{4\ttcw}<rdfs:range\ rdf:resource="http://www.w3.org/2001/XMLSchema\#string"/>}
\codeline{\hspace*{2\ttcw}</owl:DatatypeProperty>}
\codeblank{}
\codeline{\hspace*{2\ttcw}\textcolor{cmtgray}{<!--\ 实例数据\ -->}}
\codeline{\hspace*{2\ttcw}<mfg:Lathe\ rdf:about="http://example.org/manufacturing\#Lathe\_001">}
\codeline{\hspace*{4\ttcw}<rdfs:label\ xml:lang="zh">1号数控车床</rdfs:label>}
\codeline{\hspace*{4\ttcw}<mfg:serialNumber>CNC-LATHE-2024-001</mfg:serialNumber>}
\codeline{\hspace*{2\ttcw}</mfg:Lathe>}
\codeblank{}
\codeline{</rdf:RDF>}
\end{codebox}
\color{black}

\assetsource{ch04}{rdf-xml-examples}

把它与上一小节的 Turtle 并排看，关键是认出\textbf{同一批三元组的另一副面孔}：
\texttt{owl:Class} 元素（其 \texttt{rdf:about} 属性指明主语）对应
\texttt{mfg:Equipment a owl:Class}；
嵌套其中的 \texttt{rdfs:subClassOf} 元素（\texttt{rdf:resource} 指明宾语）
对应一条 \texttt{rdfs:subClassOf} 三元组；
结尾的 \texttt{mfg:Lathe} 元素是「类型化节点元素」写法，
一个标签同时完成「声明资源」与「断言 \texttt{rdf:type}」两件事。
\textbf{解析之后，两个文件得到的三元组集合完全一致}——格式只是外衣，图才是本体。

RDF/XML 的口碑不佳有其原因：同一张图存在多种合法 XML 写法
（属性形式与元素形式可互换、嵌套深浅随意），
嵌套结构诱导「树状思维」而 RDF 本质是图；
更要命的是通用 XML 工具几乎帮不上忙——
用 XML Schema 校验 RDF/XML 内容基本不可行，
因为「合法变体太多」恰恰是 Schema 这种逐字面校验最怕的局面。
还有一个冷知识级缺陷：RDF/XML 要求谓语 IRI 能拆成
「命名空间 + 合法 XML 元素名」，本地名以数字开头之类的 IRI
根本写不出来——同一图在 Turtle 里合法、在 RDF/XML 里无法序列化。
工程建议因此是\textbf{只读不写}：理解它以便对接遗留系统，
新文件优先用 Turtle。生态工具普遍能做格式互转和语法检查；但本书不再给出
某个外部工具的独立转换命令，也不把“能解析”当成图等价证明。若要把互转写成
本书的机器主张，必须在 ch04 包内补齐输入、标准化图比较与 exact oracle；当前
round-trip 合同缺失，所以 release 保持 \texttt{blocked}。

\subsection{N-Triples 与 JSON-LD}

另外两种格式各有专属生态位，一笔带过但场景要记牢。

\textbf{N-Triples} 是 Turtle 的「全展开」子集：不准用前缀、不准用分号逗号，
每行一条三元组、所有 IRI 写全。人读起来痛苦，机器处理却极舒服——
按行切分即可并行加载、流式处理，是大规模数据交换（dump 文件）的标准选择。
其四列扩展 \textbf{N-Quads} 在每行末尾追加命名图（named graph）标识，
用于传输多图数据集（4.6 节谈元数据管理时还会遇到命名图）。
感受一下同一条知识在 N-Triples 里的样子——一行就是一条完整三元组，
此处仅因排版折行：

\begin{codebox}
\codeline{<http://example.org/manufacturing\#Lathe\_001>}
\codeline{\hspace*{2\ttcw}<http://www.w3.org/1999/02/22-rdf-syntax-ns\#type>}
\codeline{\hspace*{2\ttcw}<http://example.org/manufacturing\#Lathe>\ .}
\codeline{<http://example.org/manufacturing\#Lathe\_001>}
\codeline{\hspace*{2\ttcw}<http://example.org/manufacturing\#power>}
\codeline{\hspace*{2\ttcw}"22.0"\^{}\^{}<http://www.w3.org/2001/XMLSchema\#decimal>\ .}
\end{codebox}

\textbf{JSON-LD} 则面向 Web 开发者：它是合法的 JSON，
靠一个 \texttt{@context} 块把普通 JSON 字段名映射到 IRI，
前端工程师无需理解 RDF 也能产出标准三元组。
搜索引擎的结构化数据标注（schema.org 生态）正是其最大用武之地，
JSON-LD 1.1 于 2020 年成为 W3C 推荐标准。
同样两条知识的 JSON-LD 形态，对前端工程师几乎零门槛：

\begin{codebox}
\codeline{\{}
\codeline{\hspace*{2\ttcw}"@context":\ \{\ "mfg":\ "http://example.org/manufacturing\#",}
\codeline{\hspace*{16\ttcw}"power":\ \{\ "@id":\ "mfg:power"\ \}\ \},}
\codeline{\hspace*{2\ttcw}"@id":\ \ \ "mfg:Lathe\_001",}
\codeline{\hspace*{2\ttcw}"@type":\ "mfg:Lathe",}
\codeline{\hspace*{2\ttcw}"power":\ 22.0}
\codeline{\}}
\end{codebox}

四种格式的取舍可以浓缩成一张表：

\begin{center}
\small
\begin{tabularx}{\linewidth}{@{}>{\ttfamily}p{0.14\linewidth}p{0.26\linewidth}X@{}}
\toprule
\normalfont\textbf{格式} & \textbf{一句话定位} & \textbf{适用场景} \\
\midrule
Turtle & 人读人写的事实标准 & 手写本体、教学示例、代码评审与 git 版本管理 \\
RDF/XML & 最早的标准格式 & 对接遗留系统与老工具；OWL 时代存量文件的主要形态 \\
N-Triples & 一行一条、全部展开 & 大规模数据交换与备份、流式与并行加载、命令行处理 \\
JSON-LD & JSON 的 RDF 化身 & Web API 与前端集成、搜索引擎结构化数据（schema.org） \\
\bottomrule
\end{tabularx}
\end{center}

选择口诀：\textbf{手写用 Turtle，交换用 N-Triples，上 Web 用 JSON-LD，
遇到遗留系统再谈 RDF/XML}。无论选哪种，记住它们承载的是同一张图——
任何合格的 RDF 工具都能在格式之间无损往返。
「同图异文」甚至延伸到了 HTTP 层：Linked Data 实践中，
同一个资源 IRI 可以按请求头的内容协商（content negotiation）
分别返回 Turtle、JSON-LD 或人读的 HTML 页面——
名字只有一个，表现形式随消费者而变，这正是整套设计哲学的缩影。

\section{RDFS：给事实加上第一层模式}

裸 RDF 只有事实没有词汇表：谁都可以发明谓语，机器却不知道这些谓语意味着什么。
RDFS（RDF Schema）补上第一层模式能力——它本身也用三元组书写，
核心词汇一只手数得过来，但每一个都自带\textbf{推理语义}。

\subsection{类与类层次：subClassOf}

\texttt{rdfs:Class} 声明一个资源是类，\texttt{rdfs:subClassOf} 声明类间的包含关系——
「所有车床都是数控机床」写作：
\texttt{mfg:Lathe rdfs:subClassOf mfg:CNCMachine}。两条性质要点：
\textbf{传递性}——子类的子类还是子类，层次因此可以任意深；
\textbf{多重继承合法}——一个类可以同时是多个类的子类
（「数控车床」既在「车床」之下又在「数控设备」之下），
RDFS 对此毫无限制，是否使用多继承是建模风格问题（第 3 章的方法论已有讨论）。

\texttt{rdfs:subPropertyOf} 的推理行为与类层次一致：
断言「\texttt{Spindle\_01} 直接部件于 \texttt{Lathe\_001}」，
若「直接部件于」是「部件于」的子属性，
推理器自动补出「\texttt{Spindle\_01} 部件于 \texttt{Lathe\_001}」——
数据只写最具体的关系，泛化层面交给推理，这个分工后文模式一节大量使用。

随手可得的还有注释词汇 \texttt{rdfs:label} 与 \texttt{rdfs:comment}：
它们对推理没有任何影响，对\textbf{人}却是头等重要——
标签是界面与查询结果里的展示名，注释是建模意图的随身文档。
第 3 章强调的命名规范与文档化要求，落到语言层就是坚持给每个类写这两条。

\begin{noteblock}
\textbf{rdfs:Class 还是 owl:Class？}仓库的 RDFS 教学资产用前者，
OWL 文件用后者。在工程实践里可以这样把握：
读到 \texttt{rdfs:Class}，通常意味着文件停留在纯 RDFS 表达力层；
新建本体时直接写 \texttt{owl:Class}——Protégé 等工具也只会产出后者。
理解前者是为了读懂旧文件与轻量词表，不是为了在新项目里使用它。
\end{noteblock}

\subsection{属性与 domain/range}

RDFS 把属性当一等公民：\texttt{rdf:Property} 声明属性，
\texttt{rdfs:subPropertyOf} 构造属性层次
（「直接部件于」是「部件于」的子属性，4.6 节的部分-整体模式会用到），
\texttt{rdfs:domain} 与 \texttt{rdfs:range} 则声明属性两端的类型：

\begin{itemize}
  \item \texttt{mfg:canProcess rdfs:domain mfg:Equipment}——
        凡是\textbf{作主语}用了 \texttt{canProcess} 的资源，都是设备；
  \item \texttt{mfg:canProcess rdfs:range mfg:Material}——
        凡是\textbf{作宾语}出现在 \texttt{canProcess} 之后的资源，都是材料。
\end{itemize}

请反复读上面两句话的措辞：是「都是」，不是「必须是」。
这一字之差是 RDFS 被误解最深的地方，下一小节专门掰开揉碎。

还有一个少有人知却时常致祸的细节：同一属性声明\textbf{多个} domain
（或 range），语义是\textbf{交集}而非并集——
给 \texttt{canProcess} 同时声明 domain 为设备和刀具，
意思是每个使用者「既是设备又是刀具」，而不是「设备或刀具皆可」；
想表达「或」，要把并类（OWL 的 \texttt{or} 表达式）作成一个 domain。
许多公开本体里的诡异推理结论，源头就是无意间叠出来的多重 domain。

\subsection{RDFS 推理语义：domain 是推断，不是校验}

RDFS 的标准语义由一组\textbf{衍推规则}（entailment rules）定义——
给定一张图，规则不断推出新三元组，直到没有新结论为止。核心四条：

\begin{itemize}
  \item \textbf{类继承}（rdfs9）：\(x\) 属于类 \(A\)，且 \(A \sqsubseteq B\)
        \(\Rightarrow\) \(x\) 属于 \(B\)；类层次传递（rdfs11）与之配套。
  \item \textbf{属性继承}（rdfs7）：\(x\) 与 \(y\) 有子属性关系
        \(\Rightarrow\) 它们也有父属性关系。
  \item \textbf{定义域}（rdfs2）：三元组 \((x, p, y)\) 存在且 \(p\) 的
        domain 是 \(C\) \(\Rightarrow\) \(x\) 属于 \(C\)。
  \item \textbf{值域}（rdfs3）：三元组 \((x, p, y)\) 存在且 \(p\) 的
        range 是 \(C\) \(\Rightarrow\) \(y\) 属于 \(C\)。
\end{itemize}

规则编号出自 RDF 语义规范，记号不必背，行为必须懂——
尤其要注意每条规则都只会\textbf{产生}三元组，没有任何一条会报错或拒绝数据。

包内资产把最常用的类继承推理写成了可手算的最小案例：

\begin{codebox}
\codeline{\textcolor{cmtgray}{\#\ ==============================}}
\codeline{\textcolor{cmtgray}{\#\ RDFS推理示例}}
\codeline{\textcolor{cmtgray}{\#\ ==============================}}
\codeline{\textcolor{cmtgray}{\#\ 基于以下事实：}}
\codeline{\textcolor{cmtgray}{\#\ \ \ mfg:Lathe\ rdfs:subClassOf\ mfg:Equipment}}
\codeline{\textcolor{cmtgray}{\#\ \ \ mfg:Lathe\_001\ rdf:type\ mfg:Lathe}}
\codeline{\textcolor{cmtgray}{\#\ RDFS推理器可以推出：}}
\codeline{\textcolor{cmtgray}{\#\ \ \ mfg:Lathe\_001\ rdf:type\ mfg:Equipment\ \ (类继承推理)}}
\codeblank{}
\codeline{mfg:Lathe\ a\ rdfs:Class\ ;}
\codeline{\hspace*{4\ttcw}rdfs:subClassOf\ mfg:Equipment\ ;}
\codeline{\hspace*{4\ttcw}rdfs:label\ "车床"@zh\ .}
\codeblank{}
\codeline{mfg:Lathe\_001\ a\ mfg:Lathe\ ;}
\codeline{\hspace*{4\ttcw}rdfs:label\ "1号车床"@zh\ ;}
\codeline{\hspace*{4\ttcw}mfg:locatedIn\ mfg:WorkCell\_A\ .}
\codeline{\textcolor{cmtgray}{\#\ \(\rightarrow\)\ 通过\ rdfs:subClassOf\ 推理，可得\ mfg:Lathe\_001\ a\ mfg:Equipment}}
\codeblank{}
\codeline{mfg:WorkCell\_A\ a\ mfg:WorkCell\ ;}
\codeline{\hspace*{4\ttcw}rdfs:label\ "A区工作单元"@zh\ .}
\end{codebox}

\assetsource{ch04}{rdfs-examples}

文件里\textbf{从未写过}「\texttt{Lathe\_001} 是设备」这条三元组，
按 RDFS 蕴涵规则应当推出这条类型断言；因此一个声称支持该规则的实现，必须在
绑定夹具上让「所有 \texttt{mfg:Equipment} 实例」包含 \texttt{Lathe\_001}。
这仍是规范性预期，不是 ch04 当前场景已经证明的能力；只有加入 package runner、
单因反例与 exact oracle 后，才能成为本书的可复算收据。
这是本书反复出现的主题「显式知识之外的增值」最朴素的形态。

domain/range 的推断同样如此。给定
\texttt{mfg:producedBy rdfs:range mfg:Equipment} 与事实
「某产品 \texttt{producedBy} 某资源 X」，推理器会推出「X 是设备」——
即便 X 从未被显式声明过任何类型。这在数据干净时是便利
（半结构化来源的数据自动获得类型），在数据脏时是放大器
（错误断言长出错误类型）。全部推理结果的集合称为图的
\textbf{衍推闭包}（entailment closure），标准生态中有两种取用方式：
装载时\textbf{物化}——算好写入存储投影，查询零开销、数据更新需重算；
或查询时\textbf{动态推理}——结论随输入更新、每次查询都付推理成本。
本书不把这两种外部实现直接当证据：闭包必须在 Semantica 声明的 profile 内重算，
或作为候选输出回到 package 经过 exact oracle 与 receipt。
4.5.2 小节的 RL Profile 与 4.7.4 小节的属性路径，分别是两条路线的代表。

现在看那个高频误解。某团队给 \texttt{canProcess} 声明了
domain 为 \texttt{Equipment}，本意是「拦住把订单当设备用的脏数据」。
某天数据管道误产出 \texttt{mfg:Order\_007 mfg:canProcess mfg:Aluminum\_6061}，
他们期待推理器报错——实际发生的是：推理器\textbf{安静地推出}
\texttt{mfg:Order\_007 a mfg:Equipment}。
订单从此「是」一台设备，下游所有按设备类检索的查询都会带出这张订单。
没有报警、没有不一致，错误数据反而长出了错误分类。
更糟的是发现时点：这类污染往往潜伏到几周后的某张报表数字对不上才暴露，
回溯时还得先理解「订单怎么会被算成设备」的推理链——
事故复盘的时间成本远超当初少写一层校验省下的功夫。

\begin{noteblock}
\textbf{为什么会这样}：RDFS（以及 OWL）的语义是\textbf{开放世界}下的\textbf{推断}逻辑——
公理描述世界「是什么样」，新信息只会触发新结论，不会触发「违规」。
RDFS 甚至没有任何表达「不」的手段，连不一致都制造不出来。
想要\textbf{校验}语义（缺了报错、类型不符报错），用第 7 章的 SHACL，
它在封闭检查的立场上产出违规报告；在 OWL 中声明类不相交（如订单与设备
\texttt{owl:disjointWith}）也能让上述数据触发不一致——但那已是 OWL 的能力。
「拿 domain 当校验」是新手在生产事故里学会的第一课，希望你在这里提前学会。
\end{noteblock}

\subsection{RDFS 的表达力边界}

RDFS 能干的事到此为止。盘点它\textbf{说不出口}的话：
不能说「设备与材料不相交」（无否定与不相交），
不能说「每台设备至多有一个序列号」（无基数约束），
不能说「合格产品\textbf{当且仅当}通过检验的产品」（无等价定义，只有单向包含），
不能说「先于关系可传递」「位于关系是函数性的」（无属性特性），
也不能把状态值限定在四个枚举个体之内（无枚举类），
甚至不能断言「这两个名字指同一台设备」或「这两个名字一定不同」——
个体同一性词汇（\texttt{owl:sameAs}、\texttt{owl:differentFrom}）也属于 OWL。
一句话总结：\textbf{RDFS 只能往上走（沿层次泛化），不能说「不」，也不能数数}。
工程约束恰恰大量需要后两者——这正是 OWL 登场的理由。

顺带划清一条容易混淆的界线：\texttt{rdfs:label}、\texttt{rdfs:comment}
以及 SKOS 的各类标注词汇都属于\textbf{注释性词汇}，
在 OWL 中归为注解属性（annotation property），对逻辑推理完全透明——
删光所有标签，推理结论一条不变。
区分「进逻辑的公理」与「给人看的注解」，
是阅读任何本体文件时要戴上的第一副眼镜；
也因此，4.4.1 小节才敢说「标签注释对人头等重要」而不担心它干扰语义。

\section{OWL 2 全景}

OWL（Web Ontology Language）把第 2 章的描述逻辑装进了 Web 标准的外壳。
本节按「全貌--子语言--语法--构造器--属性公理--经典陷阱」的顺序展开，
它是全章篇幅最重的一节，因为这里的每个决定都直接进入生产系统。

\subsection{OWL 与描述逻辑的对应}

OWL 2 有两种读法。\textbf{OWL 2 DL} 读法把本体看作描述逻辑
\(\mathcal{SROIQ}(D)\) 的公理集（呼应第 2 章的 DL 家族命名法：
\(\mathcal{S}\) 是带传递属性的 \(\mathcal{ALC}\)，
\(\mathcal{R}\) 加属性链与属性层次，\(\mathcal{O}\) 加枚举（nominal），
\(\mathcal{I}\) 加逆属性，\(\mathcal{Q}\) 加限定基数，\((D)\) 加数据类型）——
表达力强且推理\textbf{可判定}，是工程默认选择；
\textbf{OWL 2 Full} 读法保持与 RDF 完全兼容、不设任何语法限制
（类可以同时是个体，元建模随意），代价是推理不可判定，只在特殊场合使用。
本书此后说 OWL 均指 OWL 2 DL。

OWL 2 DL 为可判定性付出的语法代价里，有两条新手常撞上：
其一，类、属性、个体的名字默认分属不同空间，
同名复用要走受控的「双关」机制（punning）且有限制——
对象属性与数据属性之间则绝对不许同名；
其二，传递属性这类「非简单属性」不能再参与基数约束。
Protégé 会在你越界时直接报错，背后的理由可回看第 2 章的可判定性讨论。

还有一件「同图异文」的升级版事实：OWL 本体的抽象结构与具体语法分离，
同一本体可写成函数式语法（Functional Syntax，规范用）、
RDF 序列化（Turtle/RDF/XML，交换用）与 Manchester 语法（人读人写，编辑器用）。
4.5.3 小节精讲 Manchester，因为它是 Protégé 界面的语言。

工程上还有两件常被忽略的语言设施值得点名。
\textbf{模块化导入}：\texttt{owl:imports} 让一个本体声明式地引入另一个本体，
团队据此把大本体拆成「核心概念 + 各领域扩展」的模块结构，
Protégé 的本体导入面板管理这条机制（仓库 README 的导入步骤即为此）。
\textbf{版本标识}：\texttt{owl:versionIRI} 给本体的每个发布版本一个独立 IRI，
配合变更日志注解，下游系统才能明确声明「我兼容哪一版」。

\subsection{三个 Profile：工程化的表达力裁剪}

完整 OWL 2 DL 的推理复杂度极高（一致性检验为 N2ExpTime 完全），
本体一大或数据一多就可能算不动。「指数级」不是修辞——
它意味着推理时间可能随公理规模发生雪崩式增长，
同一推理器面对规模相近、结构略异的两个本体，耗时可以相差几个数量级。
标准无法让最坏情况消失，但可以圈出「保证算得快」的安全区，
这就是 Profile 的全部哲学：\textbf{把表达力关进笼子，换取复杂度的承诺}。
OWL 2 标准据此预制了三个
\textbf{子语言}（Profile）：各自砍掉一部分表达力，换取可规模化的推理算法。
\textbf{选 Profile 是本体项目最重要的工程决策之一}，
它决定了你能写哪些公理、能用哪些推理器、能扛多大数据。

\begin{itemize}
  \item \textbf{OWL 2 EL}——为「概念巨多」的场景设计，
        源自 \(\mathcal{EL}^{++}\) 逻辑家族。保留存在限制（\(\exists\)）与交（\(\sqcap\)），
        放弃全称限制、否定、析取与逆属性，换来\textbf{多项式时间}的本体分类。
        生物医学大本体是它的主场：数十万概念的 SNOMED~CT 临床术语系统、
        Gene Ontology 都按 EL 设计，ELK 推理器可在秒级完成分类。
  \item \textbf{OWL 2 QL}——为「数据躺在关系库」的场景设计，源自 DL-Lite 家族。
        表达力剪裁到恰好能把概念查询\textbf{重写成 SQL}
        （查询回答的数据复杂度落在 \(\mathrm{AC}^0\)），
        于是本体可以作为语义视图直接架在现有关系数据库之上，数据不动、语义加层——
        这套架构有个专名：\textbf{本体数据访问}（Ontology-Based Data Access, OBDA）。
  \item \textbf{OWL 2 RL}——为「实例巨多」的场景设计，
        是能用\textbf{前向规则}完整实现的最大子集。
        标准生态中的三元组库常把它实现为装载时物化（GraphDB、Jena 规则引擎等），
        以查询时零额外开销换取更新重算成本。这里是在解释历史实现策略；
        本书只承认 Semantica runtime 在声明 profile 内重算并绑定到 receipt 的结论，
        外部物化闭包只能作为待重验的候选投影。
\end{itemize}

\begin{center}
\small
\begin{tabularx}{\linewidth}{@{}p{0.09\linewidth}p{0.15\linewidth}p{0.22\linewidth}p{0.21\linewidth}X@{}}
\toprule
\textbf{Profile} & \textbf{理论渊源} & \textbf{换来的计算性质} & \textbf{放弃的表达力（举要）} & \textbf{典型场景} \\
\midrule
EL & \(\mathcal{EL}^{++}\) 家族 & 多项式时间分类 & 全称限制、否定、析取、逆属性 & 超大概念层次（SNOMED CT、基因本体） \\
QL & DL-Lite 家族 & 查询重写为 SQL，数据复杂度 \(\mathrm{AC}^0\) & 传递属性、较强基数约束 & 关系库上的语义访问（OBDA） \\
RL & 规则可实现子集 & 前向物化、装载时算完 & 推出「新个体」的存在推断等 & 大规模实例数据的图谱推理 \\
完整 DL & \(\mathcal{SROIQ}(D)\) & 可判定但最坏指数级以上 & —— & 概念与数据规模都可控的精细建模 \\
\bottomrule
\end{tabularx}
\end{center}

选型口诀：\textbf{TBox 巨大选 EL，ABox 巨大选 RL，数据在关系库里选 QL，
两头都不大才放心用完整 DL}。三点补充。
其一，Profile 是\textbf{上限承诺}而非强制——
按 EL 设计的本体放进完整 DL 推理器毫无问题，反向则要删公理，
所以宁可一开始保守。
其二，「落在某 Profile 内」是可以\textbf{机器判定}的：
W3C 在线校验器与各类工具都能检查一份本体违反了哪条 Profile 限制，
把这一检查挂进持续集成，能防止某次提交无意间把全队拖出安全区——
比如某位同事顺手加了一条 \texttt{only} 公理，EL 资格即刻作废。
其三，制造业知识系统常见的输入形态是
「中等 TBox + 大 ABox + 关系库共存」，遗留实践往往把 RL 物化与 QL 风格的
数据库映射并用。本书用这一现象解释 profile 取舍，但第 9 章只保留
Semantica package/runtime/receipt 一条执行链，外部映射与物化均是 adapter 候选。

\subsection{Manchester 语法精讲}

机器交换本体用 RDF 序列化即可，\textbf{人}写公理需要更友好的皮肤。
Manchester 语法用英语关键字替代逻辑符号，是 Protégé 类表达式编辑器的语言，
也是团队评审公理时投影在大屏上的语言。先看仓库的类定义骨架：

\begin{codebox}
\codeline{Class:\ :Equipment}
\codeline{\hspace*{4\ttcw}Annotations:}
\codeline{\hspace*{8\ttcw}rdfs:label\ "设备"@zh,}
\codeline{\hspace*{8\ttcw}rdfs:label\ "Equipment"@en,}
\codeline{\hspace*{8\ttcw}rdfs:comment\ "制造系统中所有物理设备的基类"@zh}
\codeline{\hspace*{4\ttcw}SubClassOf:\ owl:Thing}
\codeline{\hspace*{4\ttcw}DisjointWith:\ :Material,\ :Tool}
\codeblank{}
\codeline{Class:\ :ProcessingEquipment}
\codeline{\hspace*{4\ttcw}Annotations:}
\codeline{\hspace*{8\ttcw}rdfs:label\ "加工设备"@zh}
\codeline{\hspace*{4\ttcw}SubClassOf:\ :Equipment}
\codeblank{}
\codeline{Class:\ :CNCMachine}
\codeline{\hspace*{4\ttcw}Annotations:}
\codeline{\hspace*{8\ttcw}rdfs:label\ "数控机床"@zh}
\codeline{\hspace*{4\ttcw}SubClassOf:\ :ProcessingEquipment}
\codeline{\hspace*{4\ttcw}SubClassOf:\ :hasControl\ value\ :CNCControl}
\codeblank{}
\codeline{Class:\ :Lathe}
\codeline{\hspace*{4\ttcw}Annotations:}
\codeline{\hspace*{8\ttcw}rdfs:label\ "车床"@zh}
\codeline{\hspace*{4\ttcw}SubClassOf:\ :CNCMachine}
\end{codebox}
\color{black}

\assetsource{ch04}{owl-classes}

语法骨架一目了然：\texttt{Class:} 开块，缩进的小节各管一类公理——
\texttt{Annotations:} 挂标签注释，\texttt{SubClassOf:} 声明父类或必要条件，
\texttt{DisjointWith:} 声明不相交，（后文还会见到）\texttt{EquivalentTo:} 声明等价定义。
对应的属性侧关键字是 \texttt{ObjectProperty:} 与 \texttt{DataProperty:} 开块，
\texttt{Domain:}、\texttt{Range:}、\texttt{Characteristics:} 分述两端类型与特性。
类表达式内部则用小写连接词：\texttt{and}、\texttt{or}、\texttt{not}、
\texttt{some}、\texttt{only}、\texttt{value}、\texttt{min}、\texttt{max}、
\texttt{exactly}。同一本体若导出为函数式语法或 Turtle，
这些关键字会变成 \texttt{ObjectSomeValuesFrom} 或
\texttt{owl:someValuesFrom} 之类的形态——名字更长，结构相同，
评审公理认 Manchester、交换文件认 RDF 序列化，是团队里合理的分工。

三处值得驻足。第一，\texttt{Equipment} 上的
\texttt{DisjointWith: :Material, :Tool} 正是 RDFS 给不出的「说不」能力——
有了它，任何把设备同时断言为材料的数据都会让推理器报\textbf{不一致}，
第 1 章所说的「公理即约束」落到代码就是这一行。
第二，\texttt{CNCMachine} 的
\texttt{SubClassOf: :hasControl value :CNCControl} 是\textbf{值限制}：
每台数控机床都带着指向特定个体 \texttt{:CNCControl} 的关系。
第三，\texttt{Lathe} 的 \texttt{SubClassOf: :canProcess some :RotationalPart}
是\textbf{存在限制}：是车床就必然能加工某种回转体零件。

还有一个贯穿建模生涯的区分：\textbf{SubClassOf 给必要条件，
EquivalentTo 给充要条件}。「车床 \texttt{SubClassOf} 能加工回转体」
只说车床都如此，推理器\textbf{不会}把能加工回转体的东西自动归为车床；
而后文 \texttt{QualifiedProduct} 用 \texttt{EquivalentTo} 定义，
推理器才会把满足条件的个体\textbf{自动分类}进来。
想让推理器替你「认出」成员，必须用等价定义——这是 OWL 自动分类的开关。
Manchester 语法的应用场景也不止编辑器：
Protégé 的 DL Query 即席查询、推理解释面板展示肇事公理，
用的都是这套语法——它是工程师与本体之间的「工作语言」，
值得练到能脱稿读写常见公理的程度。

\subsection{类表达式构造器一览}

OWL 的类表达式像乐高：原子类是基本砖块，构造器把它们拼成任意复杂的匿名类。
下表把 Manchester 写法、描述逻辑记号与读法对齐（\(C, D\) 为类表达式，
\(R\) 为对象属性，\(a_i\) 为个体，\(n\) 为非负整数）：

\begin{center}
\small
\begin{tabularx}{\linewidth}{@{}>{\ttfamily}p{0.24\linewidth}p{0.13\linewidth}X@{}}
\toprule
\normalfont\textbf{Manchester 写法} & \textbf{DL 记号} & \textbf{读法与示例} \\
\midrule
C and D & \(C \sqcap D\) & 交：既是设备又是高功率的东西 \\
C or D & \(C \sqcup D\) & 并：车床或铣床 \\
not C & \(\lnot C\) & 补：不是测量设备的东西 \\
R some C & \(\exists R.C\) & 存在限制：至少加工一种铝合金材料 \\
R only C & \(\forall R.C\) & 全称限制：所加工的只能是金属材料 \\
R value a & \(\exists R.\{a\}\) & 值限制：控制系统是那一台 \texttt{CNCControl} \\
R min n C & \(\ge n\,R.C\) & 至少 \(n\) 个：至少配两把刀具 \\
R max n C & \(\le n\,R.C\) & 至多 \(n\) 个：至多占一个工位 \\
R exactly n C & \(= n\,R.C\) & 恰好 \(n\) 个：恰好一个序列号 \\
\{a1, a2, ...\} & \(\{a_1, a_2, \dots\}\) & 枚举类：状态只能取那几个个体 \\
\bottomrule
\end{tabularx}
\end{center}

构造器可以任意嵌套，工程上却要克制：嵌套越深，公理越难评审、推理越慢。
经验法则是\textbf{给复杂表达式起名}——把嵌套子表达式提成具名类再引用，
本体因此长出更多「中间概念」，这通常是好事而非坏事。
另一个进阶事实顺带知晓：公理两侧其实都允许放\textbf{类表达式}
（所谓一般概念包含，GCI），
「凡是能加工钛合金的设备都得有冷却系统」可以直接写成
「(\texttt{Equipment and canProcess some Titanium})
\texttt{SubClassOf} (\texttt{hasCooling some CoolingSystem})」——
左边是匿名类，无须先起名。Protégé 在类视图外提供了
「General class axioms」入口承载这类公理。

关于\textbf{枚举类}（enumeration）多说一句，它是受控值集的逻辑化身：
包内资产 \texttt{owl-classes} 用
\texttt{EquivalentTo: \{ :Idle, :Running, :Maintenance, :Faulty \}}
把设备状态封死在四个个体之内。但要记住 OWL \textbf{没有唯一名称假设}
（Unique Name Assumption, UNA）：不同 IRI 默认\textbf{可能}指同一个体，
所以枚举类务必配合 \texttt{owl:AllDifferent} 声明个体两两互异，
否则「至多一种状态」之类的基数推理会得出离谱结论。4.6.4 小节回到这个模式。

\subsection{属性公理：特性即推理行为}

类表达式之外，OWL 的另一半力量在属性公理上。仓库的属性文件值得整段精读：

\begin{codebox}
\codeline{\textcolor{cmtgray}{\#\ ==============================}}
\codeline{\textcolor{cmtgray}{\#\ 对象属性\ (Object\ Properties)}}
\codeline{\textcolor{cmtgray}{\#\ ==============================}}
\codeblank{}
\codeline{ObjectProperty:\ :canProcess}
\codeline{\hspace*{4\ttcw}Annotations:}
\codeline{\hspace*{8\ttcw}rdfs:label\ "能加工"@zh}
\codeline{\hspace*{4\ttcw}Domain:\ :Equipment}
\codeline{\hspace*{4\ttcw}Range:\ :Material}
\codeblank{}
\codeline{ObjectProperty:\ :hasCapability}
\codeline{\hspace*{4\ttcw}Annotations:}
\codeline{\hspace*{8\ttcw}rdfs:label\ "具有能力"@zh}
\codeline{\hspace*{4\ttcw}Domain:\ :Equipment}
\codeline{\hspace*{4\ttcw}Range:\ :EquipmentCapability}
\codeline{\hspace*{4\ttcw}Characteristics:\ Functional}
\codeblank{}
\codeline{ObjectProperty:\ :hasStatus}
\codeline{\hspace*{4\ttcw}Annotations:}
\codeline{\hspace*{8\ttcw}rdfs:label\ "具有状态"@zh}
\codeline{\hspace*{4\ttcw}Domain:\ :Equipment}
\codeline{\hspace*{4\ttcw}Range:\ :EquipmentStatus}
\codeline{\hspace*{4\ttcw}Characteristics:\ Functional}
\codeblank{}
\codeline{ObjectProperty:\ :locatedIn}
\codeline{\hspace*{4\ttcw}Annotations:}
\codeline{\hspace*{8\ttcw}rdfs:label\ "位于"@zh}
\codeline{\hspace*{4\ttcw}Domain:\ :Equipment}
\codeline{\hspace*{4\ttcw}Range:\ :WorkCell}
\codeblank{}
\codeline{ObjectProperty:\ :precedes}
\codeline{\hspace*{4\ttcw}Annotations:}
\codeline{\hspace*{8\ttcw}rdfs:label\ "先于"@zh}
\codeline{\hspace*{8\ttcw}rdfs:comment\ "工序之间的先后顺序，具有传递性"@zh}
\codeline{\hspace*{4\ttcw}Domain:\ :Operation}
\codeline{\hspace*{4\ttcw}Range:\ :Operation}
\codeline{\hspace*{4\ttcw}Characteristics:\ Transitive}
\codeblank{}
\codeline{ObjectProperty:\ :requires}
\codeline{\hspace*{4\ttcw}Annotations:}
\codeline{\hspace*{8\ttcw}rdfs:label\ "需要"@zh}
\codeline{\hspace*{4\ttcw}Domain:\ :Operation}
\codeline{\hspace*{4\ttcw}Range:\ :Equipment}
\codeblank{}
\codeline{ObjectProperty:\ :assignedTo}
\codeline{\hspace*{4\ttcw}Annotations:}
\codeline{\hspace*{8\ttcw}rdfs:label\ "分配给"@zh}
\codeline{\hspace*{4\ttcw}Domain:\ :Task}
\codeline{\hspace*{4\ttcw}Range:\ :Equipment}
\codeblank{}
\end{codebox}
\color{black}

\assetsource{ch04}{owl-properties}

\texttt{Characteristics:} 一行字，改变的是推理器的行为方式：

\begin{itemize}
  \item \textbf{函数性}（Functional）：每个主语至多一个值。
        \texttt{hasStatus} 声明为函数性后，若数据里出现
        「\texttt{Lathe\_001} 状态为 \texttt{Idle}」与「状态为 \texttt{Running}」两条，
        推理器\textbf{不是报错}，而是先推出 \texttt{Idle owl:sameAs Running}
        ——又是无 UNA 在起作用；只有当二者已被声明互异（如枚举类配套的
        \texttt{AllDifferent}）时才报不一致。理解这个两段式行为，
        许多「推理器怎么不报错」的困惑会瞬间消解。
  \item \textbf{传递性}（Transitive）：\texttt{precedes} 声明传递后，
        工序 A 先于 B、B 先于 C，推理器自动补出 A 先于 C——
        工艺路线的先后约束因此只需维护相邻工序对。
  \item \textbf{对称性}（Symmetric）与\textbf{逆属性}（InverseOf）：
        「相邻」声明对称则一条数据顶两条；
        把 \texttt{producedBy} 声明为 \texttt{produces} 的逆，
        两个方向的查询就共享同一份数据。
  \item \textbf{属性链}（OWL 2 新增）：
        \texttt{locatedIn o isPartOf SubPropertyOf: locatedIn} 式的公理表达
        「位于某单元、单元属于某车间，则位于该车间」——
        跨一跳的常识从此不必靠应用代码拼接。
  \item \textbf{其余特性一笔}：自反（Reflexive）、非自反（Irreflexive）、
        非对称（Asymmetric）补全了 OWL 2 的特性清单；
        注意\textbf{数据属性只能声明函数性}——
        其余特性都以「个体对个体」的关系为前提。
  \item \textbf{键与逆函数性}：逆函数性（InverseFunctional）声明
        「值唯一确定主语」，相当于对象属性版的全局键；
        数据属性想充当键（设备按序列号全局唯一），
        用 OWL 2 的 \texttt{owl:hasKey} 公理表达——
        两台序列号相同的「设备」会被推理器合并成同一个体。
\end{itemize}

\begin{noteblock}
\textbf{特性选错是一类隐蔽缺陷}。把 \texttt{locatedIn} 误声明为函数性，
设备搬迁后保留的历史位置记录会触发「两个工作单元是同一个」的荒谬推理；
把本不传递的「直接部件」声明为传递，BOM 的层级语义即刻被冲垮（见 4.6.3）。
公理不是文档注释，\textbf{每一条特性声明都会被推理器当真}——
上线前用推理器跑一遍一致性，再用第 3 章的 CQ 回归验证行为，缺一不可。
\end{noteblock}

\subsection{存在与全称：辨析与「全称陷阱」}

\(\exists\)（some）与 \(\forall\)（only）是新手最易混用的一对构造器，
也是仓库专门准备了对照说明文件的地方。先看两个真实定义：

\begin{codebox}
\codeline{\textcolor{cmtgray}{\#\ ==============================}}
\codeline{\textcolor{cmtgray}{\#\ OWL约束（属性限制）}}
\codeline{\textcolor{cmtgray}{\#\ ==============================}}
\codeline{\textcolor{cmtgray}{\#\ OWL\ Property\ Restrictions}}
\codeblank{}
\codeline{\textcolor{cmtgray}{\#\ 存在限制（\(\exists\)）：每个产品至少经过一次检验}}
\codeline{\textcolor{cmtgray}{\#\ ExistentialRestriction:\ hasInspection\ some\ Inspection}}
\codeblank{}
\codeline{Class:\ :QualifiedProduct}
\codeline{\hspace*{4\ttcw}Annotations:}
\codeline{\hspace*{8\ttcw}rdfs:label\ "合格产品"@zh}
\codeline{\hspace*{4\ttcw}EquivalentTo:}
\codeline{\hspace*{8\ttcw}:Product}
\codeline{\hspace*{8\ttcw}and\ (:hasInspection\ some}
\codeline{\hspace*{12\ttcw}(:Inspection}
\codeline{\hspace*{13\ttcw}and\ (:hasResult\ value\ "Pass")))}
\codeblank{}
\codeline{\textcolor{cmtgray}{\#\ 全称限制（\(\forall\)）：高端设备的所有特性都是高级特性}}
\codeline{\textcolor{cmtgray}{\#\ UniversalRestriction:\ hasFeature\ only\ AdvancedFeature}}
\codeblank{}
\codeline{Class:\ :HighEndEquipment}
\codeline{\hspace*{4\ttcw}Annotations:}
\codeline{\hspace*{8\ttcw}rdfs:label\ "高端设备"@zh}
\codeline{\hspace*{4\ttcw}EquivalentTo:}
\codeline{\hspace*{8\ttcw}:Equipment}
\codeline{\hspace*{8\ttcw}and\ (:price\ some\ xsd:decimal[>=\ "1000000"\^{}\^{}xsd:decimal])}
\codeline{\hspace*{8\ttcw}and\ (:hasFeature\ only\ :AdvancedFeature)}
\end{codebox}
\color{black}

\assetsource{ch04}{owl-classes}

\texttt{QualifiedProduct} 展示了\textbf{嵌套的存在限制}——
合格产品就是：产品，且\textbf{存在}一次检验，
该检验又\textbf{存在}一个取值为 Pass 的结果。
写成第 2 章的 DL 模板是 \(\exists R.(C \sqcap \exists S.\{v\})\)，
其中 \(R\)＝hasInspection、\(C\)＝Inspection、\(S\)＝hasResult、\(v\)＝Pass。
\texttt{HighEndEquipment} 则把存在与全称合用：
价格条款用 \texttt{some}（存在一个不低于百万的价格），
特性条款用 \texttt{only}（所有特性都得是高级特性）。
价格条款里还藏着一个值得学会的写法：
\texttt{xsd:decimal[>= "1000000"\^{}\^{}xsd:decimal]} 是\textbf{数据范围刻面}
（facet）——直接在公理里限定数值区间，
「功率不低于 15 千瓦」「利用率介于 0 与 1 之间」都靠它表达。
存在与全称两者的语义差异，包内教学资产总结得很精准：

\begin{codebox}
\codeline{\textcolor{cmtgray}{\#\ ==============================}}
\codeline{\textcolor{cmtgray}{\#\ 3.\ 两者的区别}}
\codeline{\textcolor{cmtgray}{\#\ ==============================}}
\codeline{\textcolor{cmtgray}{\#\ \(\exists\)R.C（存在限制）：至少一个，即"存在"}}
\codeline{\textcolor{cmtgray}{\#\ \ \ \(\rightarrow\)\ 如果满足，推理器保证存在一个相应个体}}
\codeline{\textcolor{cmtgray}{\#\ \ \ \(\rightarrow\)\ 适合表达"必须有..."}}
\codeline{\textcolor{cmtgray}{\#}}
\codeline{\textcolor{cmtgray}{\#\ \(\forall\)R.C（全称限制）：全部满足，即"所有"}}
\codeline{\textcolor{cmtgray}{\#\ \ \ \(\rightarrow\)\ 如果有关系存在，则所有关系对象都必须满足约束}}
\codeline{\textcolor{cmtgray}{\#\ \ \ \(\rightarrow\)\ 如果没有关系，则全称限制自动满足（空满足）}}
\codeline{\textcolor{cmtgray}{\#\ \ \ \(\rightarrow\)\ 适合表达"只允许..."}}
\codeblank{}
\end{codebox}

\assetsource{ch04}{property-restrictions}

重点盯住全称限制的第二行注释：\textbf{如果没有关系，全称限制自动满足}——
逻辑上称为\textbf{空满足}（vacuous satisfaction）。
这就是著名的\textbf{全称陷阱}：假如 \texttt{HighEndEquipment}
的定义里只有 \texttt{hasFeature only :AdvancedFeature} 一个条款，
那么一台\textbf{没有任何特性记录}的破旧设备会被推理器堂堂正正地分类为高端设备——
它的「所有」特性（零个）确实都是高级特性。
仓库定义靠 \texttt{price some ...} 条款挡住了这个漏洞：
\texttt{some} 强制存在性，\texttt{only} 才有的放矢。
\textbf{「only 不蕴含 some」}，要表达「有且仅有」，两者必须同时写——
这个组合在建模里常见到有个昵称叫「闭合公理」（closure axiom）。

闭合公理的具体写法值得见一次真容。「车床只加工回转体与盘类零件」写作：
\texttt{canProcess some RotationalPart} \textbf{and}
\texttt{canProcess only (RotationalPart or DiscPart)}——
\texttt{some} 保证有活可干，\texttt{only} 圈定能干的范围，
括号里的析取列出全部许可项。漏写 \texttt{some} 是空满足漏洞，
漏写 \texttt{only} 则范围形同虚设，成对出现才是完整语义。

开放世界假设让 \texttt{only} 的杀伤力进一步放大：
推理器永远假设「你没说的事可能存在」，
因此 \texttt{only} 很难被「证实」，却很容易在与其它公理互动时制造意外不一致。
包内资产 \texttt{property-restrictions} 第 4 节给出的工程建议值得抄在手边：
\textbf{优先用 \(\exists\)，谨慎用 \(\forall\)，复杂表达式拆开写，少嵌套}。

\section{高级建模模式}

语言给了砖头，但承重墙怎么砌是另一门手艺。
本节收录四个生产系统里出镜率最高的建模模式，
每个都按「问题--方案--Turtle 示意」展开。
它们大多源自 W3C 工作组备忘录与本体设计模式（Ontology Design Patterns, ODP）社区，
属于经过千锤百炼的公共财产——遇到对应问题时请直接复用，不要重新发明。
读模式时建议带着第 3 章的眼光：每个模式本质上都是
「某类能力问题反复出现」沉淀下来的标准答案，
识别出「我的问题属于哪个模式」往往比写公理本身更见功力。

\subsection{n 元关系：用中介节点拆高维事实}

\textbf{问题}：三元组天生只能表达\textbf{二元}关系，
而真实业务满地都是高维事实——
「\texttt{Lathe\_001} 在 2024 年 3 月 1 日 10 点测得主轴温度 41.5 摄氏度」
牵涉设备、时间、测点、数值四个参与者，硬塞进一条三元组无处下脚。

\textbf{方案}：把这件\textbf{事}本身提升为一个节点（W3C 备忘录称为中介节点，
实践中常按「事件/观测/记录」命名），让每个参与者用一条二元关系挂上去：

\begin{codebox}
\codeline{\textcolor{cmtgray}{\#\ 「设备--时间--测点--数值」四元事实\ \(\rightarrow\)\ 一个观测节点\ +\ 四条二元边}}
\codeline{mfg:Reading\_20240301\_0001\ a\ mfg:SensorReading\ ;}
\codeline{\hspace*{4\ttcw}mfg:ofEquipment\ \ \ \ \ \ \ mfg:Lathe\_001\ ;}
\codeline{\hspace*{4\ttcw}mfg:measuredProperty\ \ mfg:SpindleTemperature\ ;}
\codeline{\hspace*{4\ttcw}mfg:hasValue\ \ \ \ \ \ \ \ \ \ "41.5"\^{}\^{}xsd:decimal\ ;}
\codeline{\hspace*{4\ttcw}mfg:atTime\ \ \ \ \ \ \ \ \ \ \ \ "2024-03-01T10:00:00"\^{}\^{}xsd:dateTime\ .}
\end{codebox}

代价是查询多绕一跳（先找观测节点再取值），换来的是任意维度可扩展——
明天要加「操作员」「量具」字段，再挂两条边即可。
中介节点模式还有个隐形好处：观测自身成了一等公民，
质量元数据（量具编号、校准状态、置信度）从此都有安放之处——
在追溯审计严格的行业，这不是锦上添花，而是合规刚需。
给中介节点起\textbf{具名 IRI}（如带日期流水号）而不是用空节点，
理由见 4.2.4 小节。传感器领域已有现成本体走这条路线
（W3C 的 SOSA/SSN 观测本体），自研之前值得先看看。
顺带的反模式预警：有人试图把时间塞进属性名
（\texttt{hasTemperature\_20240301}）、或把多个值拼进一个字符串
（\texttt{"41.5@2024-03-01"}）——前者让属性无限增殖、模式失控，
后者杀死全部数值查询能力，两者都会在第一个统计需求面前原形毕露。

\subsection{物化：对断言本身做断言}

\textbf{问题}：有时需要谈论\textbf{三元组本身}——
「\texttt{Lathe\_001} 位于 \texttt{WorkCell\_A}」这条记录是谁录入的？何时？置信度多少？
三元组没法直接当另一条三元组的主语。

\textbf{方案一}是 RDF 标准自带的\textbf{物化}（reification）四件套：

\begin{codebox}
\codeline{\textcolor{cmtgray}{\#\ 把一条三元组「实体化」成\ rdf:Statement\ 节点，再挂元数据}}
\codeline{mfg:stmt\_042\ a\ rdf:Statement\ ;}
\codeline{\hspace*{4\ttcw}rdf:subject\ \ \ mfg:Lathe\_001\ ;}
\codeline{\hspace*{4\ttcw}rdf:predicate\ mfg:locatedIn\ ;}
\codeline{\hspace*{4\ttcw}rdf:object\ \ \ \ mfg:WorkCell\_A\ ;}
\codeline{\hspace*{4\ttcw}mfg:assertedBy\ mfg:Engineer\_Zhang\ ;}
\codeline{\hspace*{4\ttcw}mfg:assertedAt\ "2024-02-18"\^{}\^{}xsd:date\ .}
\end{codebox}

它能用，但两个缺点出了名：每条断言膨胀成四条三元组起步；
更隐蔽的是 \texttt{rdf:Statement} 节点\textbf{并不自动断言原三元组}——
上面这段写完，「\texttt{Lathe\_001} 位于 \texttt{WorkCell\_A}」本身仍需另写一条，
两者在语义上是脱钩的。

\textbf{方案二}是\textbf{命名图}（named graph）：把一批三元组装进一个带 IRI 的图，
元数据（来源、批次、入库时间）挂在图 IRI 上，
N-Quads/TriG 格式与所有主流三元组库都原生支持——
按「数据源」「批次」粒度管理出处时，这是首选。
\textbf{方案三}是标准化进行中的 RDF-star（RDF 1.2 方向）：
用 \texttt{<<\ s\ p\ o\ >>} 语法直接引用三元组，单条断言级的注记最优雅，
主流库已陆续支持，采用前确认你的存储与工具链版本。
经验法则：\textbf{批量出处用命名图，单条注记用 RDF-star（或退回 n 元事件建模），
经典物化只在对接老系统时使用}。
还要分清两个同名概念：本小节的「物化」（reification）是\textbf{建模}动作——
把断言变成可谈论的节点；4.4.3 小节的「物化」（materialization）
是\textbf{推理}动作——把衍推结论算出来存盘。
中文恰好撞名，英文与语境能把它们分开。

\subsection{部分--整体：partOf 的传递性设计}

\textbf{问题}：制造业的物料清单（BOM）是天然的部分--整体层级：
轴承装在主轴里，主轴装在车床上。业务要双向问：
「\texttt{Bearing\_07} 最终属于哪台设备？」（穿透查询）
与「主轴的\textbf{直接}子件有哪些？」（单层清单）。
只建一个传递属性，穿透有了、单层没了；只建非传递属性则相反。

\textbf{方案}：双属性模式——非传递的「直接部件」作子属性，
传递的「部件」作父属性，断言只写直接层，穿透交给推理：

\begin{codebox}
\codeline{\textcolor{cmtgray}{\#\ 双属性模式：直接部件（不传递）\ \(\sqsubseteq\)\ 部件（传递）}}
\codeline{mfg:isPartOf\ a\ owl:ObjectProperty,\ owl:TransitiveProperty\ .}
\codeline{mfg:isDirectPartOf\ a\ owl:ObjectProperty\ ;}
\codeline{\hspace*{4\ttcw}rdfs:subPropertyOf\ mfg:isPartOf\ .}
\codeblank{}
\codeline{\textcolor{cmtgray}{\#\ 数据只维护直接层级}}
\codeline{mfg:Bearing\_07\ mfg:isDirectPartOf\ mfg:Spindle\_01\ .}
\codeline{mfg:Spindle\_01\ mfg:isDirectPartOf\ mfg:Lathe\_001\ .}
\codeline{\textcolor{cmtgray}{\#\ 推理自动补出：mfg:Bearing\_07\ mfg:isPartOf\ mfg:Lathe\_001}}
\end{codebox}

两个配套提醒。第一，\textbf{partOf 不是 subClassOf}：
「主轴是车床的部件」（个体组成关系）与「车床是机床的子类」（概念包含关系）
是两种完全不同的语义，混用是第 3 章 OntoClean 体检常抓的典型病——
判别口诀是看「实例会不会跟着走」：车床的实例都是机床的实例，
但主轴的实例绝不是车床的实例。
第二，部分--整体其实有多种细分（组件/成员/材料构成等），
W3C 的简单部分--整体关系备忘录给出了细致分类，建复杂产品结构前值得一读。
查询侧的配合见 4.7.4 小节：若存储未开传递推理，
\texttt{mfg:isPartOf+} 属性路径同样能完成穿透——
模式设计与查询习语是一对，建模时就该想好查询会怎么写。

\subsection{受控词表与枚举值集}

\textbf{问题}：状态、等级、类别这类「取值有限」的字段若放任为自由字符串，
\texttt{"idle"}、\texttt{"IDLE"}、\texttt{"空闲"} 三个写法很快各自为政，
统计与过滤全部失真。

\textbf{方案 A}：取值封闭且很小（个位数）时，用 4.5.4 小节的\textbf{枚举类}
把值集写进逻辑，让推理器替你守门：

\begin{codebox}
\codeline{\textcolor{cmtgray}{\#\ 枚举值集：状态类的外延被\ owl:oneOf\ 封死为四个个体}}
\codeline{mfg:EquipmentStatus\ a\ owl:Class\ ;}
\codeline{\hspace*{4\ttcw}owl:oneOf\ (\ mfg:Idle\ mfg:Running\ mfg:Maintenance\ mfg:Faulty\ )\ .}
\codeblank{}
\codeline{mfg:Idle\ a\ mfg:EquipmentStatus\ ;\ rdfs:label\ "待机"@zh\ .}
\codeline{\textcolor{cmtgray}{\#\ 别忘了\ owl:AllDifferent\ 声明四个个体两两互异（无\ UNA！）}}
\end{codebox}

\textbf{方案 B}：词表大、会演化、还有层次（材料牌号上千、工艺分类逐年修订）时，
不要把它压进类层次，改用 \textbf{SKOS}（Simple Knowledge
Organization System）受控词表：每个值是一个 \texttt{skos:Concept}，
\texttt{skos:prefLabel} 与 \texttt{skos:altLabel} 管首选名与别名，
\texttt{skos:broader} 组织宽窄层次，
整表归入一个 \texttt{skos:ConceptScheme}。
SKOS 概念是\textbf{个体}而非类，增删值不动 TBox、不触发重新分类，
治理成本低得多。选择标准一句话：\textbf{进逻辑用 oneOf，进治理用 SKOS}——
判断的试金石是问自己：「推理器需要拿这些值参与公理运算吗？」
需要（如基数约束、不相交推理）就走枚举类；
只是分类、检索、对齐用，SKOS 的轻量与灵活完胜。

\section{SPARQL 深入}

SPARQL 之于 RDF，就像 SQL 之于关系表；但它的心智模型不是表连接，
而是\textbf{图模式匹配}：写一个带变量的「图模板」，
引擎在数据图中找出所有能让模板成立的变量绑定。
本节从基本形态一路推进到优化与安全。书中查询的唯一复算入口是 ch04/ch09 package；
只有绑定具名场景、输入与 oracle 的查询才可宣称通过，CQ7 的外部
\texttt{SERVICE} 在离线运行中必须 fail closed。

\subsection{基本图模式与 SELECT}

一条 SELECT 查询四段式：前缀声明、\texttt{SELECT} 投影变量、
\texttt{WHERE} 图模式、尾部修饰（排序分页）。
\texttt{WHERE} 里的三元组模式与 Turtle 语法几乎一致（分号逗号照用），
只是任何位置都可以放 \texttt{?变量}；
多条模式以\textbf{合取}组成\textbf{基本图模式}（Basic Graph Pattern, BGP），
共享变量就是天然的连接条件。第 3 章说「每条能力问题翻译成一条 SPARQL 即验收标准」，
现在看 CQ1 的标准答案：

\begin{codebox}
\codeline{\textcolor{cmtgray}{\#\ CQ1:\ 哪些设备可以加工铝合金？}}
\codeline{\textcolor{cmtgray}{\#\ Which\ equipment\ can\ process\ aluminum\ alloy?}}
\codeline{\textcolor{cmtgray}{\#\ ==============================}}
\codeline{SELECT\ ?equipment\ ?name\ ?power}
\codeline{WHERE\ \{}
\codeline{\hspace*{4\ttcw}?equipment\ a\ :Equipment\ ;}
\codeline{\hspace*{15\ttcw}:canProcess\ :Aluminum\ ;}
\codeline{\hspace*{15\ttcw}rdfs:label\ ?name\ .}
\codeline{\hspace*{4\ttcw}OPTIONAL\ \{\ ?equipment\ :power\ ?power\ \}}
\codeline{\hspace*{4\ttcw}FILTER(lang(?name)\ =\ "zh")}
\codeline{\}}
\codeline{ORDER\ BY\ ASC(?name)}
\end{codebox}

\assetsource{ch04}{sparql-queries}

逐行拆解：\texttt{?equipment a :Equipment} 限定类型；
\texttt{:canProcess :Aluminum} 用\textbf{常量}收紧候选集；
\texttt{rdfs:label ?name} 取展示名。三条模式共享 \texttt{?equipment}，
引擎对每个候选设备同时满足三条才放行——这就是图模式的合取语义。
查询的产出是一张\textbf{绑定表}：每行是一组「变量到值」的映射，
\texttt{SELECT} 挑选要投影的列；
尾部修饰符与 SQL 同款——\texttt{ORDER BY} 排序、
\texttt{LIMIT} 与 \texttt{OFFSET} 分页，
\texttt{DISTINCT} 去重在多路径可达同一答案时尤其常用。

\begin{noteblock}
\textbf{为什么不直接用 SQL？}关系库里这条 CQ 至少要在设备表、能力表、
标签表之间写两三个 JOIN，而且「沿类层次找所有子类设备」在 SQL
里要么递归公用表表达式、要么应用层循环。
图模式把「沿边导航」做成语言原语，层次穿透交给属性路径（4.7.4），
这是查询语言与数据模型匹配带来的红利。
\end{noteblock}

\subsection{OPTIONAL、UNION 与 FILTER}

\texttt{OPTIONAL} 是 SPARQL 版的左外连接：花括号内的模式\textbf{匹配则绑定、
不匹配不淘汰整行}。CQ1 里功率字段就是这样处理的——
没有功率记录的设备也要出现在结果里，只是 \texttt{?power} 列空着。
若写成普通模式，缺数据的设备会整行消失，
报表「莫名少了几台设备」的事故多半源于此。仓库 CQ8 是更立体的示范：

\begin{codebox}
\codeline{\textcolor{cmtgray}{\#\ ==============================}}
\codeline{\textcolor{cmtgray}{\#\ CQ8:\ 查询故障设备及其维修历史}}
\codeline{\textcolor{cmtgray}{\#\ ==============================}}
\codeline{SELECT\ ?equipment\ ?faultTime\ ?maintenanceRecord}
\codeline{WHERE\ \{}
\codeline{\hspace*{4\ttcw}?equipment\ a\ :Equipment\ ;}
\codeline{\hspace*{15\ttcw}:hasStatus\ :Faulty\ ;}
\codeline{\hspace*{15\ttcw}:hasFaultTime\ ?faultTime\ .}
\codeline{\hspace*{4\ttcw}OPTIONAL\ \{}
\codeline{\hspace*{8\ttcw}?maintenanceRecord\ :relatedTo\ ?equipment\ ;}
\codeline{\hspace*{27\ttcw}a\ :MaintenanceRecord\ .}
\codeline{\hspace*{4\ttcw}\}}
\codeline{\}}
\codeline{ORDER\ BY\ DESC(?faultTime)}
\end{codebox}

\assetsource{ch04}{sparql-queries}

故障设备必须有故障时间（必选模式），维修记录则可有可无（OPTIONAL 块）——
新设备首次故障时还没有维修历史，但它必须出现在排查清单上。
这种「核心字段必选、外围字段可选」的分层写法是生产查询的基本素养；
顺带留意 OPTIONAL 块内部可以是\textbf{多条}模式（此处两条），
它们作为整体要么全部匹配、要么整块落空，
拆成两个 OPTIONAL 与合在一个里语义并不相同。

\texttt{UNION} 表达\textbf{或}：两个花括号块各自匹配，结果取并集，
常用于「同义异构」的数据——例如同时兼容新旧两版属性命名的过渡期查询。
\texttt{FILTER} 则对已绑定的变量施加布尔条件：
数值比较（\texttt{?power > 20}）、语言筛选（\texttt{lang(?name) = "zh"}）、
正则匹配（\texttt{regex(?sn, "\^{}CNC-")}）。
还有一个常用否定习语：\texttt{FILTER NOT EXISTS \{ ... \}}
——「不存在某模式」，下文 ASK 一节用它写约束校验。
否定还有一位近亲 \texttt{MINUS}：从左侧结果中减去与右侧匹配的行。
它与 \texttt{FILTER NOT EXISTS} 在多数场景结果一致，
但当两侧不共享任何变量时行为大相径庭（此时 \texttt{MINUS} 什么都不减）——
团队里最好约定只用其中一个，省下排查诡异差异的整个下午。

\subsection{聚合与子查询}

SPARQL 1.1 把 SQL 的统计能力搬了过来：
\texttt{GROUP BY} 分组；\texttt{COUNT}、\texttt{SUM}、\texttt{AVG}、
\texttt{MIN}、\texttt{MAX}、\texttt{GROUP\_CONCAT} 聚合；
\texttt{HAVING} 对聚合值过滤：

\begin{codebox}
\codeline{\textcolor{cmtgray}{\#\ 各工作单元的设备数与平均功率，只看设备数\ >=\ 2\ 的单元}}
\codeline{SELECT\ ?workCell\ (COUNT(?eq)\ AS\ ?cnt)\ (AVG(?p)\ AS\ ?avgPower)}
\codeline{WHERE\ \{}
\codeline{\hspace*{4\ttcw}?eq\ a\ :Equipment\ ;}
\codeline{\hspace*{8\ttcw}:locatedIn\ ?workCell\ ;}
\codeline{\hspace*{8\ttcw}:power\ ?p\ .}
\codeline{\}}
\codeline{GROUP\ BY\ ?workCell}
\codeline{HAVING\ (COUNT(?eq)\ >=\ 2)}
\codeline{ORDER\ BY\ DESC(?cnt)}
\end{codebox}

规则与 SQL 同源：\texttt{SELECT} 里出现的非聚合变量必须进 \texttt{GROUP BY}；
想「先算后比」就用\textbf{子查询}——内层 \texttt{SELECT} 先求出
每个单元的平均功率，外层再挑出功率高于本单元平均值的设备。
子查询自内向外执行、只通过投影变量与外层连接，
是把复杂统计拆成可读分层的主要工具。
仓库 CQ4（统计各单元设备数）是练习聚合的现成起点。

再记一个小而常用的子句 \texttt{VALUES}：把一组常量内联成绑定表。
例如\linebreak[4]\texttt{VALUES ?st \{ :Maintenance :Faulty \}}
一行就把查询限定在「维护中或故障」两种状态上——
比一长串 \texttt{FILTER} 的或条件清爽得多，
也是参数化查询模板的天然挂点（第 8 章自动生成查询时会再见到它）。

\subsection{属性路径：subClassOf* 详讲}

\textbf{属性路径}（property path）让谓语位置可以写正则风格的表达式：
\texttt{/} 串接，\texttt{|} 择一，\texttt{\^{}} 反向，
\texttt{*}（零或多步）、\texttt{+}（一或多步）、\texttt{?}（零或一步）。
其中最重要的习语，怎么强调都不过分：

\begin{codebox}
\codeline{\textcolor{cmtgray}{\#\ 标准习语：取某类「连同其全部子类」的实例}}
\codeline{SELECT\ ?eq\ WHERE\ \{}
\codeline{\hspace*{4\ttcw}?eq\ rdf:type/rdfs:subClassOf*\ :Equipment\ .}
\codeline{\}}
\codeblank{}
\codeline{\textcolor{cmtgray}{\#\ 传递可达：工序\ Op\_10\ 之后（直接或间接）的全部工序}}
\codeline{SELECT\ ?later\ WHERE\ \{}
\codeline{\hspace*{4\ttcw}:Op\_10\ :precedes+\ ?later\ .}
\codeline{\}}
\end{codebox}

第一条要会默写：\texttt{?eq} 的类型沿 \texttt{rdfs:subClassOf}
走\textbf{零步或多步}到达 \texttt{:Equipment}——
于是车床、数控机床的实例统统命中，哪怕端点\textbf{没开任何推理}。
这招常被称为「穷人版推理」：用路径在查询时现场爬类层次。
但要清醒：它只补了类层次这一种衍推，
domain/range、属性层次、传递属性带来的结论一概不补；
\texttt{*}/\texttt{+} 在深层次大图上也可能代价高昂。
在标准生态里，图谱够大、衍推种类多时常见做法是 RL 物化（4.5.2）；
轻量场景与即席分析则常用属性路径。在本书执行边界内，两者都只是候选策略：
必须由 Semantica 对绑定输入重算，通过 query/shape/oracle 回归并写入 receipt，
才可支撑正式主张。

反向符 \texttt{\^{}} 的妙用值得单独一提：
\texttt{:Aluminum\_6061 \^{}:canProcess ?eq} 从材料出发反查设备，
与 \texttt{?eq :canProcess :Aluminum\_6061} 完全等价——
当导航需要在长路径中间「掉头」时（例如沿部件层级向上、再沿位置关系向下），
反向符能省掉一串中间变量，让一条路径写完整个导航语义。

\subsection{CONSTRUCT 与 ASK}

\texttt{SELECT} 之外还有三个查询形态。\texttt{CONSTRUCT} 做\textbf{图到图}变换：
\texttt{WHERE} 匹配出的每组绑定按模板生成新三元组，
常用于抽取子图、做轻量 ETL、把高频推理结论物化成显式候选数据。
候选图若要进入本书语义边界，仍须回到 Semantica 绑定、重验证并生成 receipt。
\texttt{ASK} 只回答真假，是把约束类 CQ 变成\textbf{自动回归测试}的利器
（第 3 章说过：这类检查期望结果为空，非空即告警）：

\begin{codebox}
\codeline{\textcolor{cmtgray}{\#\ CONSTRUCT：把「设备\ \(\rightarrow\)\ 所在单元中文名」固化为新图}}
\codeline{CONSTRUCT\ \{\ ?eq\ :inCellNamed\ ?label\ \}}
\codeline{WHERE\ \{\ ?eq\ :locatedIn\ ?c\ .\ ?c\ rdfs:label\ ?label\ .}
\codeline{\hspace*{8\ttcw}FILTER(lang(?label)\ =\ "zh")\ \}}
\codeblank{}
\codeline{\textcolor{cmtgray}{\#\ ASK：还有没有位置的设备吗？期望\ false，true\ 即数据告警}}
\codeline{ASK\ \{\ ?eq\ a\ :Equipment\ .}
\codeline{\hspace*{6\ttcw}FILTER\ NOT\ EXISTS\ \{\ ?eq\ :locatedIn\ ?cell\ \}\ \}}
\end{codebox}

第四个形态 \texttt{DESCRIBE} 返回某资源的「相关描述」，
内容由端点自行决定，探索陌生数据集时顺手，正式接口里少用。
把 \texttt{ASK} 与 \texttt{SELECT} 组合进持续集成是成熟团队的标配：
CQ 清单存成 \texttt{.sparql} 文件、与本体一起进版本库，
每次提交自动全量回放、比对期望结果——
第 3 章倡导的「CQ 即验收」从口号变成流水线，靠的正是这套机制。

\subsection{查询优化直觉}

执行计划是引擎的事，但写查询的人有三条朴素直觉能避免大多数性能事故：

\begin{itemize}
  \item \textbf{高选择性模式往前放}：带常量（具体 IRI、具体值）的模式先把
        候选集打小，纯变量模式（\texttt{?s ?p ?o} 之流）千万别打头阵。
  \item \textbf{先收窄再外延}：\texttt{OPTIONAL} 与代价高的 \texttt{FILTER}
        放在核心模式之后；现代引擎会自动下推过滤，但清晰的书写顺序
        同时服务于人与机器。
  \item \textbf{警惕笛卡尔积与无界路径}：不共享变量的模式组合会产生
        绑定爆炸；深层 \texttt{*} 路径在大图上同样危险。
        调试期养成 \texttt{LIMIT} 护身的习惯。
\end{itemize}

拿 CQ1 体会第一条：先匹配 \texttt{:canProcess :Aluminum}
（能加工铝合金的设备屈指可数），还是先匹配 \texttt{?equipment a :Equipment}
（全厂设备成百上千），中间结果规模差着量级；
现代优化器多半会替你重排，但在解释计划、排查慢查询时，
这套「选择性」直觉就是你的地图。
一句话总纲：\textbf{让引擎尽早拿到小集合}。复杂报表查询拆成子查询分层聚合，
比一口气写两屏的「巨无霸」更快也更可维护。

\subsection{SPARQL UPDATE 的历史语义与受控 mutation 边界}

SPARQL 1.1 的更新语言在标准生态中可以把端点从只读升级为可写：
\texttt{INSERT DATA} 插入既定三元组（仓库 CQ5 新增 3 号车床即此形态），
\texttt{DELETE ... INSERT ... WHERE} 按模式批量改写
（仓库 CQ6 把已派工的空闲设备批量置为运行中），
\texttt{LOAD}/\texttt{CLEAR}/\texttt{DROP} 管理整图。
这些语句值得学习，因为它们精确表达图级 mutation；但在本书现行架构中，
外部 SPARQL 端点只是 Semantica 已验收 Dataset 的\textbf{只读投影}，不提供第二个
写入面。所有 mutation 都先形成候选 delta，交回有权限的 Semantica 执行面完成
内容绑定、shape/query/oracle 重验证并生成新 receipt。标准生态中的三条防线可据此重述：

\begin{itemize}
  \item \textbf{投影只读}：Fuseki 等服务器虽能分别配置查询与更新服务，
        本书部署 adapter 只暴露查询；禁用其更新服务，而不是把它藏到内网后继续旁路写入。
  \item \textbf{操作白名单}：面向应用（尤其是由大语言模型生成查询）的投影通道
        只放行 \texttt{SELECT}/\texttt{ASK}；\texttt{INSERT}/\texttt{DELETE}/
        \texttt{LOAD}/\texttt{CLEAR}/\texttt{DROP} 一律拒绝并转成候选 delta。
        第 8 章的 Text2SPARQL 闸门设计中，「操作白名单」正是四道校验之一。
  \item \textbf{资源限额}：为公开端点配置超时与结果行数上限，
        防止深路径或笛卡尔积查询拖垮服务。
\end{itemize}

两件相邻话题顺带交代。其一，\textbf{联邦查询}（federated query）：
\texttt{SERVICE <端点>} 子句可以把局部模式发到远端端点执行
（仓库 CQ7 演示了跨端点取设备的外部编号），
跨组织协作不必先把数据搬到一处——但远端的延迟与可用性从此进入你的故障域，
作为 projection adapter 使用时务必配超时与降级路径，且远端结果要连同 snapshot、
query、adapter 身份与输出 hash 回到 Semantica receipt。其二，\textbf{注入防护}：
用字符串拼接把用户输入嵌进查询文本，与 SQL 时代的注入漏洞一模一样；
遗留 adapter 一律走参数化接口（Jena 的 \texttt{ParameterizedSparqlString}、
rdflib 的 \texttt{initBindings} 等）或用 \texttt{VALUES} 传入绑定；
参数化只防注入，不能替代 Semantica 的内容绑定与语义验收。

\section{工具链：从编辑器到推理器}

语言之后是生态比较。本节介绍交互式编辑器、编程库与推理器，帮助读者理解标准为何
可由不同工具实现；这些工具不构成本书的第二条执行链。本书相关的查询、校验、推理、
版本与 receipt 仍只经 Semantica package/runner。

\subsection{历史互操作观察：Protégé 工作流}

Protégé（斯坦福大学出品，免费开源）是常用本体开发 IDE。
下面保留一条经典建模工作流，目的是读懂既有项目与标准生态，而不是为本书建立
第二个写入或验收入口。若用 Protégé 观察 package 中的导出副本，任何编辑结果都只能
作为迁移候选；它必须回到 ch04 package，重新计算内容绑定，经过 Semantica 校验、
场景 oracle 与 receipt 后，才可能成为本书可执行语义的一部分：

\begin{enumerate}
  \item \textbf{观察本体身份}：历史工作流会在 File \(\rightarrow\) New 后于 Active Ontology
        视图设置本体 IRI（对应 4.2 节的命名问题——IRI 起好就别再改）；
        也可 File \(\rightarrow\) Open 查看只读的 \texttt{manufacturing} 导出副本。
  \item \textbf{观察类层次}：Entities 标签的 Classes 面板曾用于逐级 Add subclass，
        每建一个类立即补 \texttt{rdfs:label}（中英双语）与 \texttt{comment}——
        界面上的 Annotations 区就是 4.4.1 小节那两个词汇。
  \item \textbf{观察属性公理}：Object properties 面板可呈现对象属性、
        设 Domain/Range，在 Characteristics 区勾选函数性、传递性等——
        勾选框背后正是 4.5.5 小节的属性公理，勾之前想清楚推理后果。
  \item \textbf{观察类表达式}：类表达式编辑器用
        Manchester 语法添加 \texttt{SubClass\allowbreak Of} 或
        \texttt{Equivalent\allowbreak To} 公理，
        编辑器带自动补全与语法检查（4.5.3 小节的语法在这里日常使用）。
  \item \textbf{比较外部推断}：历史上常从 Reasoner 菜单选择 HermiT 并启动。
        推断出的层次以高亮底色显示——对照 4.4.3 小节，
        你能在界面上直接看到 \texttt{Lathe\_001} 被「搬进」\texttt{Equipment}；
        若本体不一致，问题类标红，Explain 按钮给出肇事公理集
        （称为 justification，第 5 章细讲其原理）。这类显示只是比较证据，
        不能替代 Semantica 的受限推理结果与 receipt。
  \item \textbf{比较 DL Query}：DL Query 标签可输入 Manchester 表达式
        （如 \texttt{Equipment and canProcess some Aluminum}），
        勾选 Instances 与 Subclasses 查看结果；正式验收仍须把同一问题表达为
        package 中绑定的 query/CQ，并由 Semantica 重算。
  \item \textbf{识别插件产物}：可按工具自身文档了解 OntoGraf（图形化浏览层次）、
        SWRLTab（第 5 章的规则编辑）、SPARQL Query（库内直接跑查询）。
        团队协作场景另有 WebProtégé；其导出内容同样只是待迁移输入，不自动取得权威身份。
\end{enumerate}

这条历史工作流留下的有效经验是：公理增量要小、检查要频繁，不一致越早暴露越好定位。
本版把这个节奏落在 package 版本、迁移记录和 Semantica 场景回归上，而不是把 IDE 的
绿色界面当成验收结论。

两个互操作习惯仍值得保留。\textbf{交换副本优先 Turtle}：
Protégé 另存时可能给出 RDF/XML，Turtle 通常更适合差异评审（理由见 4.3.1 小节），
但序列化文件并不因此成为 package 正本。
\textbf{分清断言层与推断层}：外部推理器开启时，
界面把手写公理与推理结论混排展示（后者有高亮标记），
导出时必须记录层次与引擎；任何显式闭包都要作为新候选资产回到 Semantica，
重新绑定输入、运行时、输出和 provenance，不能直接交付为本书结论。

\subsection{历史编程生态：四个库如何分工}

下表说明既有系统为何会出现不同库，也帮助设计迁移适配器；它不是本书的当前选型表：

\begin{center}
\small
\begin{tabularx}{\linewidth}{@{}p{0.15\linewidth}p{0.09\linewidth}p{0.30\linewidth}X@{}}
\toprule
\textbf{库} & \textbf{语言} & \textbf{推理支持} & \textbf{定位与适用} \\
\midrule
Apache Jena & Java & 内置 RDFS 与通用规则引擎，可挂接外部 OWL 推理器 & 既有 Java 系统常见：ARQ、TDB2 与 Fuseki 组成查询/存储栈 \\
Eclipse RDF4J & Java & RDFS 级推理，第三方扩展接入 & 仓库抽象层与 SDK，GraphDB 等商用库的开放接口 \\
rdflib & Python & 自身不带推理，配合 \texttt{owlrl} 包做 RDFS/OWL RL 物化 & 脚本、数据管道、快速原型，Python 数据生态的黏合剂 \\
owlready2 & Python & 捆绑调用 HermiT/Pellet & 「本体当对象」编程：类即 Python 类，改完同步回 OWL \\
\bottomrule
\end{tabularx}
\end{center}

这张表只用于读懂生态差异与迁移来源。既有项目可因语言、存储、许可与 profile 使用
Jena、RDF4J、RDFLib 或 owlready2，但进入两卷书边界后，它们只能位于受控 adapter 后面；
其输出必须重新进入 Semantica package 做内容绑定、校验和场景复算。第 9 章不再维持
Java/Jena 与 Python/owlready2 双栈：标准化红利体现在资产可审查、可迁移，而不是允许
多个引擎各自给出无法绑定的答案。

表格之外还有两个名字会经常遇到：\textbf{OWL API} 是 Java
世界操作 OWL 公理级结构的底层库，Protégé 与多数推理器都建在它之上，
直接使用它的场景多为开发本体工具而非业务应用；
Python 侧的 \textbf{SPARQLWrapper} 是访问远程 SPARQL 端点的轻量客户端。
这些名字用于识别遗留接口；本书不把它们暴露为绕过 Semantica 的可执行入口。

本书的最小组合拳只有一条：从 allowlisted package 加载内容绑定资产，用
\texttt{SemanticRuntime} 查询/校验/受限推理，由 \texttt{SemanticPackageRunner}
比较 oracle，再把 runtime/输入/输出/PROV 绑定到 receipt。编辑器可用于人工观察，
外部存储可作为部署适配面，但都不能绕过这条证据链。

\subsection{历史推理器对比与互操作边界}

推理器是把公理变成结论的引擎，核心任务四件，名字值得记准：
\textbf{一致性检验}（本体整体有没有矛盾）、
\textbf{可满足性检验}（某个类是否注定为空——空类多半是建模事故）、
\textbf{分类}（算出完整的蕴含类层次，即推断版的 subClassOf 全图）、
\textbf{实例化}（判定个体属于哪些类）；
外加一项体验型服务\textbf{解释}（给出某条结论的最小肇事公理集）。
四个主流选手各有性格：

\begin{center}
\small
\begin{tabularx}{\linewidth}{@{}p{0.13\linewidth}p{0.30\linewidth}p{0.21\linewidth}X@{}}
\toprule
\textbf{推理器} & \textbf{核心算法} & \textbf{覆盖范围} & \textbf{许可} \\
\midrule
HermiT & 超表（hypertableau）演算，削减经典 Tableau 的分支爆炸 & OWL 2 DL & LGPL \\
Pellet & 经典 Tableau，支持 SWRL 规则 & OWL 2 DL & AGPL，另提供商业授权 \\
ELK & 基于结论的饱和推理，多核并行 & OWL 2 EL（子集） & Apache 2.0 \\
Konclude & 并行 Tableau 与饱和技术结合 & OWL 2 DL & LGPL v3 \\
\bottomrule
\end{tabularx}
\end{center}

这张表应按算法能力来读，而不是照表为本书挑后端：ELK 的 EL 饱和、HermiT/Pellet/
Konclude 的完整 DL 路径，解释了 profile、完备性、增量分类、性能和许可之间的取舍。
它也解释了为何两个外部引擎可能对同一输入给出不同覆盖范围或解释结构。

若迁移项目必须与遗留引擎做对照实验，实验须固定引擎版本、profile、输入快照与配置，
产出只作为 provenance 完整的比较证据；不得把外部分类结果直接写入正式资产。
候选结论必须回到 Semantica 受支持的推理边界重新计算，经过 query/shape/oracle 回归并
生成 receipt。外部引擎无法在该边界内复现的结论，应被标记为差异或能力缺口，而不是
悄然成为第二种真值。Tableau 算法如何工作、何时终止，仍是第 5 章的理论主线。

\section{本章小结与练习}

\subsection{要点回顾}

本章信息密度不小，收尾把骨架重新拎一遍——
下列每一条，都值得在合上书后自问能否用自己的话复述：

\begin{itemize}
  \item 语义网技术栈是\textbf{分层的开放标准族}：RDF 管事实、RDFS/OWL 管模式与逻辑、
        SPARQL 管查询、SHACL 管校验；分层意味着可渐进采纳，标准化意味着厂商可替换。
  \item RDF 把知识拆成\textbf{三元组}，节点只有 IRI、字面量、空节点三种身份；
        语言标签与数据类型是字面量的质量关键，空节点「能不用就不用」。
  \item 序列化格式只是外衣：\textbf{同一张图可写成 Turtle、RDF/XML、
        N-Triples、JSON-LD}，按场景选择，随时无损互转。
  \item RDFS 的 domain/range 是\textbf{推断}而非校验——「拿 domain 当校验」
        是必须提前接种的疫苗；校验语义属于 SHACL（第 7 章）。
  \item OWL 2 提供完整描述逻辑表达力，三个 Profile（EL/QL/RL）
        是\textbf{表达力换规模}的工程套餐；SubClassOf 给必要条件，
        EquivalentTo 才触发自动分类；only 有空满足陷阱，记得闭合公理。
  \item n 元关系、物化、部分--整体、受控词表四个\textbf{建模模式}
        覆盖了高维事实、断言元数据、层级穿透、值集治理四类高频需求。
  \item SPARQL 的心智模型是\textbf{图模式匹配}；OPTIONAL 防缺字段丢行，
        「类型沿子类层次上溯」的属性路径习语要会默写，
        ASK 把约束 CQ 变成回归测试；外部 projection 端点只读，
        mutation 必须回到 Semantica 形成新绑定与 receipt。
  \item 标准生态提供多种编辑器、编程库与推理器，说明 RDF/OWL 资产具有可迁移性；
        本书的执行组合却只有 Semantica package、runtime、runner 与 receipt，
        外部工具结果不能替代这条绑定证据链。
\end{itemize}

\subsection{思考题}

\begin{enumerate}
  \item 把「2 号车床能加工钛合金 TC4」写成三元组并给出 Turtle 文本
        （自拟 IRI，注意前缀声明）；再为钛合金补充英文标签。
        追问：中文标签与英文标签是一条三元组还是两条？为什么互不覆盖？
  \item 「同图异文」实证：读取 ch04 包的 \texttt{rdf-xml-examples} 与
        \texttt{rdf-turtle-examples} 资产，比较两种序列化解析后的三元组集合。
        两个文件的三元组集合为什么可以直接比较？
        空节点（若有）会给这种比较带来什么麻烦？
  \item 给 \texttt{mfg:canProcess} 声明了 \texttt{rdfs:domain mfg:Equipment}
        之后，数据里混入 \texttt{mfg:Order\_007 mfg:canProcess mfg:Aluminum\_6061}。
        RDFS 推理器会报错吗？会发生什么？若想让这类数据被「拦下」，
        给出两种不同机制的方案并比较其立场差异。
  \item 设计题：为「设备维护事件」建模——涉及设备、维护人员、日期、
        更换的备件、费用五个参与者。用 4.6.1 小节的中介节点模式
        给出类与属性设计及一条示例数据的 Turtle 草案。
  \item 全称陷阱：若把 \texttt{HighEnd\allowbreak Equipment} 的定义削减为仅剩
        \texttt{:hasFeature only :Advanced\allowbreak Feature}，
        一台没有任何 \texttt{hasFeature} 断言的设备会被分类为高端设备吗？
        解释原因，并写出修复后的完整定义。
  \item Profile 选型：为三个场景各选一个 OWL 2 Profile 并说明理由——
        (a) 数十万概念的医疗术语层次需要每日重新分类；
        (b) 制造执行系统的数据都在 Oracle 里，业务要语义化查询但数据不许搬家；
        (c) 内容绑定 Dataset 有上亿条设备运行记录，需要按本体规则批量打标签，
        并把一个外部三元组库存储视图作为可重建投影。
  \item SPARQL 实战：写一条查询，统计每个工作单元中「数控机床及其全部子类」
        的实例数量，要求用属性路径而不开推理；再写一条 ASK 查询，
        校验「是否存在没有序列号的设备」。分别说明期望结果。
  \item 执行边界题：查看 ch04 manifest，列出哪些 ontology/query asset 会被 runtime
        加载，哪些 Manchester/\texttt{SERVICE} 能力会阻断；运行开放/封闭世界场景，
        分别记录场景 oracle 与 release verdict。解释为何不能用外部 DL 工具的输出
        替换当前 package 的 unsupported 状态。
\end{enumerate}

提示：第 3 题回看 4.4.3 小节的衍推规则与第 7 章的 SHACL 定位；
第 5 题的关键词是「空满足」（4.5.6 小节）；
第 7 题的属性路径习语在 4.7.4 小节，注意统计题要配 \texttt{GROUP BY}；
第 8 题同时使用本章章首绑定框与 4.8 节的生态比较。
个人练习产物可以保存，但若要成为本书新的可执行语义，必须以受治理 delta 进入
Semantica package，并通过来源、正反例、回归与发布门禁。

\subsection{本章包资产指引}

包资产按「数据--模式--逻辑--查询」的顺序排布，
与本章行文一一对应。建议的用法不是通读，而是\textbf{随章翻阅、动手改坏再改好}：
改一条公理看推理结论怎么变、改一个模式看查询结果怎么变，
语言的语义只有在这种反馈循环里才能长进肌肉记忆。

\begin{center}
\begin{tabularx}{\linewidth}{@{}>{\ttfamily\footnotesize}p{0.36\linewidth}X@{}}
\toprule
\normalfont\textbf{Semantica asset} & \textbf{内容与阅读建议} \\
\midrule
rdf-turtle-examples & Turtle：前缀、类层次、属性、实例，配合 4.2/4.3 节逐段读 \\
rdf-xml-examples & RDF/XML 序列化，与 Turtle 对照体会“同图异文” \\
rdfs-examples & 类层次与 domain/range；含可手算的推理案例（4.4 节） \\
owl-classes & Manchester 类定义；来源保留，当前 runtime 不解析 \\
owl-properties & Manchester 属性特性；来源保留，当前 runtime 不解析 \\
property-restrictions & \(\exists\)/\(\forall\) 语义对照教学材料 \\
sparql-queries / cq01--cq08 & 查询集合及拆分后的查询正本；并非每条都有离线 oracle \\
manufacturing & 本体教学资产；其可执行部分由 manifest 与 runtime profile 决定 \\
\bottomrule
\end{tabularx}
\end{center}

下一章我们钻进推理器内部：Tableau 算法如何从公理推出结论、
SWRL 规则如何表达业务逻辑——历史工具中的「Start reasoner」按钮背后发生了什么，
以及为何其显示结果必须回到 Semantica 受支持边界重算后才能成为本书证据。
